\documentclass[12pt,a4paper,UTF8]{ctexart}
\usepackage[a4paper,top=22mm,bottom=22mm,left=22mm,right=22mm,headheight=15pt]{geometry}
\newcommand{\Needspace}[1]{\par\begingroup\dimen0=#1\relax\dimen2=\pagegoal\advance\dimen2 by-\pagetotal\ifdim\dimen0>\dimen2\newpage\fi\endgroup}
\newcommand{\needspace}[1]{\par\begingroup\skip0=\lastskip\vskip-\skip0\dimen0=#1\relax\advance\dimen0 by\skip0\vskip\dimen0\penalty-100\vskip-\dimen0\vskip\skip0\penalty9999\endgroup}
\usepackage{amsmath,amssymb,mathtools,bm}
\usepackage{graphicx,xcolor,booktabs,longtable,array,enumitem,fancyhdr,caption,float,microtype,setspace}
\usepackage[hidelinks,unicode]{hyperref}
\usepackage[most]{tcolorbox}
\definecolor{ContestBlue}{HTML}{17365D}
\hypersetup{pdftitle={第060套},pdfauthor={}}
\setstretch{1.24}
\setlength{\parindent}{2em}
\setlength{\parskip}{0.35em}
\setlength{\emergencystretch}{2em}
\setlist{itemsep=0.35em,topsep=0.4em}
\captionsetup{font=small,labelfont=bf,justification=centering}
\setcounter{secnumdepth}{0}
\pagestyle{fancy}
\fancyhf{}
\fancyhead[L]{\small 第060套}
\fancyfoot[C]{\thepage}
\renewcommand{\headrulewidth}{0.4pt}
\ctexset{
  section={format=\Large\bfseries\color{ContestBlue},beforeskip=1.2em,afterskip=0.7em},
  subsection={format=\large\bfseries\color{ContestBlue},beforeskip=1.0em,afterskip=0.5em},
  subsubsection={format=\normalsize\bfseries,beforeskip=0.8em,afterskip=0.35em}
}
\usepackage{calc}
\newcounter{none}
\let\originalsubsection\subsection
\renewcommand{\subsection}{\needspace{6\baselineskip}\originalsubsection}
\makeatletter
\newcommand{\tightlist}{\setlength{\itemsep}{0pt}\setlength{\parskip}{0pt}}
\newcommand{\maxwidth}{\ifdim\Gin@nat@width>\linewidth\linewidth\else\Gin@nat@width\fi}
\newcommand{\maxheight}{\ifdim\Gin@nat@height>0.72\textheight 0.72\textheight\else\Gin@nat@height\fi}
\newcommand{\pandocbounded}[1]{\begingroup\setkeys{Gin}{width=\maxwidth,height=\maxheight,keepaspectratio}#1\endgroup}
\makeatother
\begin{document}
\begin{titlepage}
  \centering
  \vspace*{58mm}
  {\Huge\bfseries\color{ContestBlue} 第060套\par}
  \vfill
  \vspace*{28mm}
\end{titlepage}
\section{第 1 题}\label{ux7b2c-1-ux9898}

如图，一根细绳穿过两个质量均为 \(m\)
的小铁环。上环水平放置，绳从其左右两侧接触点 \(D,E\) 穿过，再绕下方的环
\(C\)；左右结构对称，绳与小环间的静摩擦因数为
\(\mu\)。调整上环的位置，使侧视图中绳与上方两个小环的接触点 \(D,E\)
位于同一水平面，\(AD,BE\) 与竖直方向的夹角均为 \(\alpha\)，\(DC,EC\)
近似竖直。系统保持静止。

把绳与小环的局部接触简化为绳绕圆环的问题，并以接触弧一端的张力为
\(T(0)=T_0\)，另一端张力为 \(T(\alpha)\)。

（1）若小环固定，求绳不发生相对滑动时 \(T(\alpha)\) 的最小值
\(T_{\min}\) 和最大值 \(T_{\max}\)。

（2）取 \(\mu=1.5\)，求系统能够静止时 \(\alpha\) 的取值范围。

（3）求系统存在静止平衡所需的最小静摩擦因数 \(\mu_c\)，以及相应的平衡角
\(\alpha_c\)。

\begin{figure}[H]
\centering
\includegraphics[width=0.9\linewidth,height=\textheight,keepaspectratio,alt={图1-1　对称悬挂结构与绳元在圆环接触区的受力关系}]{../figures/question-01.pdf}
\caption*{图1-1　对称悬挂结构与绳元在圆环接触区的受力关系}
\end{figure}

\subsection{一、题型定位}\label{ux4e00ux9898ux578bux5b9aux4f4d}

上下两个实体环分别平衡。上环在 \(D,E\)
接触绳，同时受自重和悬挂下环的绳段向下拉力；绳绕环的张力范围由静摩擦限制。

\subsection{二、知识点}\label{ux4e8cux77e5ux8bc6ux70b9}

\begin{enumerate}
\def\labelenumi{\arabic{enumi}.}
\tightlist
\item
  柔绳微元的法向平衡与切向平衡。
\item
  静摩擦不等式 \(|\mathrm dT|\leq\mu\,\mathrm dN\)。
\item
  绞盘公式 \(T_2/T_1=e^{\mu\varphi}\) 的推导和方向解释。
\item
  对称系统中的竖直分量平衡。
\item
  超越方程的图像法与数值求根。
\item
  临界参数对应的重根、相切与极值条件。
\end{enumerate}

\subsection{三、物理图像}\label{ux4e09ux7269ux7406ux56feux50cf}

沿接触弧取一个极短绳元。两端张力方向相差
\(\mathrm d\theta\)，它们的法向合力约为
\(T\,\mathrm d\theta\)，由小环的支持力抵消；切向张力差则由摩擦力承担。越接近即将向某一侧滑动的状态，摩擦力越接近
\(\mu\,\mathrm dN\)，张力沿弧长便按指数规律增长或衰减。

在整体装置中，下方小环由两段近似竖直的绳共同支承，故每段竖直绳张力为
\(mg/2\)。上环除自身重力外，还受到两段竖直绳向下的拉力
\(2T_0=mg\)，总共需支承 \(2mg\)，故两侧倾斜绳各提供 \(mg\)
的竖直分量，所需张力为
\(mg/\cos\alpha\)。问题由此变为：摩擦所允许的最大张力能否达到这个值。

\subsection{四、解题入口}\label{ux56dbux89e3ux9898ux5165ux53e3}

先不要处理三个小环。只看图 1-1 右侧接触弧，建立角坐标
\(\theta\)，得到张力比的上下界。再回到对称装置，用两个最简单的竖直平衡式确定
\(T_0\) 与 \(T\)。最后把 \(T\leq T_{\max}\) 写成

\[
e^{\mu\alpha}\cos\alpha\geq2.
\]

第二问是求该不等式在 \(0<\alpha<\pi/2\) 内的解集；第三问是求函数
\(g(\alpha)=e^{\mu\alpha}\cos\alpha\) 的最大值刚好等于 \(2\) 时的
\(\mu\)。

\subsection{五、解析}\label{ux4e94ux89e3ux6790}

\begin{figure}[H]
\centering
\includegraphics[width=0.487\linewidth,height=\textheight,keepaspectratio,alt={图1-2　平衡方程的两个函数交点及临界相切条件}]{../figures/solution-01.pdf}
\caption*{图1-2　平衡方程的两个函数交点及临界相切条件}
\end{figure}

图1-2中，题面图把 \(D,E\) 标在同一个上环上。由下环求 \(T_0\)
后，上环平衡给出
\(2T\cos\alpha=2mg\)；曲线交点表示张力需求落在摩擦允许范围的角度区间。

\needspace{3\baselineskip}
\textbf{（一）绳元受力与张力上下界}

取接触角为 \(\mathrm d\theta\) 的微小绳元。其一端张力为 \(T\)，另一端为
\(T+\mathrm dT\)。在法向方向，两端张力的合力保留一阶小量：

\[
\mathrm dN=T\,\mathrm d\theta. \tag{1}
\]

这里忽略了
\(\mathrm dT\,\mathrm d\theta\)，因为它是二阶小量。切向方向的张力差由静摩擦力承担，故

\[
|\mathrm dT|\leq \mu\,\mathrm dN
=\mu T\,\mathrm d\theta. \tag{2}
\]

两边除以 \(T>0\)：

\[
-\mu\,\mathrm d\theta
\leq\frac{\mathrm dT}{T}
\leq\mu\,\mathrm d\theta. \tag{3}
\]

从 \(\theta=0\) 积分到 \(\theta=\alpha\)，并使用 \(T(0)=T_0\)：

\[
-\mu\alpha
\leq\ln\frac{T(\alpha)}{T_0}
\leq\mu\alpha. \tag{4}
\]

指数函数单调递增，因此

\[
T_0e^{-\mu\alpha}
\leq T(\alpha)
\leq T_0e^{\mu\alpha}. \tag{5}
\]

于是

\[
\boxed{T_{\min}=T_0e^{-\mu\alpha}},
\qquad
\boxed{T_{\max}=T_0e^{\mu\alpha}}. \tag{6}
\]

上下界分别对应绳即将向相反方向滑动的两种极限。静止时的实际张力可以位于两界之间，不要求处处达到摩擦极限。

\needspace{3\baselineskip}
\textbf{（二）整体平衡条件}

下方小环由左右两段近似竖直、张力相同的绳支承。竖直方向平衡给出

\[
2T_0=mg,
\qquad
T_0=\frac{mg}{2}. \tag{7}
\]

上环只有一个实体，\(D,E\)
是它的两处接触点。两根倾斜绳水平分量抵消；竖直方向必须支承上环自重
\(mg\) 和下环经竖直绳传来的 \(2T_0=mg\)，因此

\[
2T\cos\alpha=2mg. \tag{8}
\]

因此所需倾斜段张力为

\[
T=\frac{mg}{\cos\alpha}. \tag{9}
\]

把式（7）、（9）代入式（5）。上界条件为

\[
\frac{mg}{\cos\alpha}
\leq\frac{mg}{2}e^{\mu\alpha}, \tag{10}
\]

即

\[
\boxed{e^{\mu\alpha}\cos\alpha\geq2}. \tag{11}
\]

下界条件为

\[
\frac{mg}{\cos\alpha}
\geq\frac{mg}{2}e^{-\mu\alpha}. \tag{12}
\]

在 \(0<\alpha<\pi/2\) 内，左边不小于 \(mg\)，右边不大于
\(mg/2\)，故式（12）自动满足。平衡范围由式（11）限制。

\needspace{3\baselineskip}
\textbf{（三）\(\mu=1.5\) 时的角度范围}

边界角满足

\[
e^{1.5\alpha}\cos\alpha=2. \tag{13}
\]

令

\[
F(\alpha)=1.5\alpha+\ln(\cos\alpha)-\ln2. \tag{14}
\]

则求根问题变为 \(F(\alpha)=0\)。由于

\[
F'(\alpha)=1.5-\tan\alpha, \tag{15}
\]

\(F\) 先增后减，在 \(\alpha=\arctan1.5\)
处取最大值，因此存在两个根。用二分法求得

\[
\alpha_1=0.5820\,\mathrm{rad},
\qquad
\alpha_2=1.2673\,\mathrm{rad}. \tag{16}
\]

式（11）在两个根之间成立，所以

\[
\boxed{0.5820\leq\alpha\leq1.2673\ \mathrm{rad}}. \tag{17}
\]

换成角度约为

\[
\boxed{33.35^\circ\leq\alpha\leq72.61^\circ}. \tag{18}
\]

\needspace{3\baselineskip}
\textbf{（四）临界摩擦因数}

临界时，曲线

\[
y_1=\cos\alpha,
\qquad
y_2=2e^{-\mu\alpha} \tag{19}
\]

只有一个交点，且在该点相切。函数值相等给出

\[
\cos\alpha_c=2e^{-\mu_c\alpha_c}. \tag{20}
\]

导数相等给出

\[
-\sin\alpha_c
=-2\mu_c e^{-\mu_c\alpha_c}. \tag{21}
\]

用式（20）消去指数项：

\[
\sin\alpha_c=\mu_c\cos\alpha_c, \tag{22}
\]

所以

\[
\boxed{\mu_c=\tan\alpha_c}. \tag{23}
\]

再代回式（20）：

\[
e^{\alpha_c\tan\alpha_c}\cos\alpha_c=2. \tag{24}
\]

数值求根得

\[
\boxed{\alpha_c=0.9144\,\mathrm{rad}=52.39^\circ}, \tag{25}
\]

\[
\boxed{\mu_c=1.2981}. \tag{26}
\]

在 \(\alpha\to0\) 时接触弧趋零，张力比趋于 \(1\)，但装置要求
\(T/T_0=2\)，所以不能平衡。在 \(\alpha\to\pi/2\)
时，所需张力发散，也不能平衡。因此可行区间位于中间角度，与结果一致。

\needspace{3\baselineskip}
\textbf{接触点与实体环的区别}

原题侧视图中的 \(D,E\) 是同一个水平上环的两边，不是两个各重 \(m\)
的环。图1-1左侧用椭圆表示这一个实体，并在 \(D,E\)
标出绳的接触位置。整体竖直方向应满足
\(2T\cos\alpha=2mg\)；对上环单独分析则有
\(2T\cos\alpha=mg+2T_0\)。这两个方程与 \(2T_0=mg\)
一致。若误画为三个各重 \(m\) 的环，则张力比会变为
\(3/\cos\alpha\)，已经改变原题模型。

推广到下环质量 \(M\) 时，\(T_0=Mg/2\)，上环受力为
\(2T\cos\alpha=(m+M)g\)，因此

\[
{T\over T_0}={m+M\over M\cos\alpha},\qquad
 e^{\mu\alpha}\cos\alpha\ge{m+M\over M}.
\]

取 \(M=m\) 恢复本题的系数 \(2\)。

\subsection{六、关键洞察}\label{ux516dux5173ux952eux6d1eux5bdf}

\begin{enumerate}
\def\labelenumi{\arabic{enumi}.}
\tightlist
\item
  绞盘公式的指数来自``每一小段摩擦能力都与当地张力成正比''，不是经验拟合。
\item
  静摩擦只能给出张力比的区间，不能在未说明即将滑动方向时直接写等号。
\item
  整体平衡把张力比固定为
  \(T/T_0=2/\cos\alpha\)，局部摩擦给出最大可实现张力比；两者比较才形成式（11）。
\item
  临界摩擦因数不是令端点角为 \(0\) 或
  \(\pi/2\)，而是让两个边界根合并为重根。
\end{enumerate}

\subsection{七、易错点}\label{ux4e03ux6613ux9519ux70b9}

\begin{enumerate}
\def\labelenumi{\arabic{enumi}.}
\tightlist
\item
  把 \(\mathrm dN\) 写成 \(2T\sin(\mathrm d\theta)\)
  却不展开，导致多出因子。
\item
  直接写
  \(T=T_0e^{\mu\alpha}\)，忽略静摩擦未必达到极限以及相反滑动趋势。
\item
  把下环的平衡误写成 \(T_0=mg\)；实际有两段绳共同支承。
\item
  第二问只求一个根，未注意函数在区间内先增后减。
\item
  临界条件只写函数值相等，未补导数相等，因而无法同时确定
  \(\alpha_c,\mu_c\)。
\end{enumerate}

\subsection{八、课堂讲解建议}\label{ux516bux8bfeux5802ux8bb2ux89e3ux5efaux8bae}

先用绳绕圆柱的演示让学生感受``包角越大、可维持的张力比越大''。板书微元时，把法向与切向分开：法向给
\(\mathrm dN\)，切向给
\(\mathrm dT\)，再明确指出二阶小量为何舍去。完成式（6）后，让学生说出正、负指数分别对应哪一侧将要滑动。

第二阶段只画 \(y=\cos\alpha\) 与
\(y=2e^{-\mu\alpha}\)。先问端点处谁大，再问两曲线是否一定相交。这样第三问的相切条件会自然出现。数值求根应给出区间和残差，例如把两根代回式（13）检查左边是否为
\(2\)。

评分时，第一问若只写绞盘公式而无微元推导，可保留结论分但不应给满；第二问必须同时写出整体张力需求和摩擦上界；第三问必须具有``相切---两方程---数值根''完整链条。

\subsection{九、变式训练}\label{ux4e5dux53d8ux5f0fux8badux7ec3}

\begin{enumerate}
\def\labelenumi{\arabic{enumi}.}
\tightlist
\item
  若下环质量为 \(M\)、上环质量仍为 \(m\)，推导平衡不等式。
\item
  若左右包角分别为
  \(\alpha_1,\alpha_2\)，摩擦因数也不同，写出两侧张力允许区间。
\item
  证明函数 \(e^{\mu\alpha}\cos\alpha\) 在 \((0,\pi/2)\)
  内至多有一个极大值。
\item
  对 \(\mu=1.5\) 用牛顿法从 \(0.5\) 和 \(1.3\) 两个初值分别求根。
\end{enumerate}

\needspace{3\baselineskip}
\textbf{变式解答：上、下环质量不等}

设上环质量为 \(M\)，下环质量为 \(m\)，仍保持对称悬挂。下环平衡给出
\(2T_0=mg\)，上环受自重和两个向下的 \(T_0\)，故

\[
2T\cos\alpha=(M+m)g,\qquad
{T\over T_0}={1+M/m\over\cos\alpha}.
\]

令 \(c=1+M/m\)，摩擦允许条件为

\[
e^{\mu\alpha}\cos\alpha\ge c.
\]

临界参数在最大值处相切，因而

\[
\mu_c=\tan\alpha_c,\qquad
\alpha_c\tan\alpha_c-\ln\sec\alpha_c=\ln c.
\]

左端对 \(\alpha\) 的导数为 \(\alpha\sec^2\alpha>0\)，所以每个 \(c>1\)
对应唯一临界角。质量比增大使 \(c\)
增大、临界角和所需摩擦同时增大。\(M=m\) 时 \(c=2\)，恢复主解析的临界
\(\mu_c\simeq1.2981\)。

这个变式还揭示``三个环''误读的影响：若把 \(D,E\)
误当两个独立上环，连质量比和受力平衡都会改变，而不只是图上多画一个圆。应先核对实体数量，再应用绕绳公式。

\needspace{3\baselineskip}
\textbf{临界附近为什么出现两个允许角端点}

记 \(f(\alpha)=\ln(2/\cos\alpha)/\alpha\)，允许条件为
\(\mu\ge f(\alpha)\)。临界点满足 \(f'(\alpha_c)=0\)。对
\(\mu=\mu_c+\Delta\mu\)，作局部展开：

\[
f(\alpha)=\mu_c+{1\over2}f''(\alpha_c)(\alpha-\alpha_c)^2+O((\alpha-\alpha_c)^3).
\]

令 \(h(\alpha)=\alpha\tan\alpha-\ln(2\sec\alpha)\)，有
\(f'=h/\alpha^2\)、\(h'=\alpha\sec^2\alpha\)，所以

\[
f''(\alpha_c)={\sec^2\alpha_c\over\alpha_c}>0.
\]

两个端点因此近似为

\[
\alpha_{\max/\min}\simeq\alpha_c
\pm\sqrt{2\alpha_c\cos^2\alpha_c\,\Delta\mu}.
\]

允许角区间宽度按 \(\sqrt{\Delta\mu}\)
增长，而不是线性增长。恰在临界摩擦时只剩一个角度；稍微增大摩擦才出现有限区间。图中两交点应表示同一允许区间的左右端点，不能当成任取两个独立的平衡状态而忽略它们之间的可行性。

近似要求 \(\Delta\mu\) 很小。原题 \(\mu=1.5\) 应用完整方程求根，得到约
\(0.5820\)、\(1.2673\) rad，不以这一局部展开替代数值结果。

\newpage

\section{第 2 题}\label{ux7b2c-2-ux9898}

两块水平木板平行放置，间距为 \(2r\)。下板固定，上板以恒定角速度
\(\Omega\bm e_z\) 绕过 \(O\) 点的竖直轴旋转。两板之间夹有一半径为 \(r\)
的球，球在上下接触点都不滑动，并且球的角速度始终没有竖直分量。以 \(O\)
为原点建立水平坐标系 \(xOy\)，球心初始坐标为 \((l,0)\)。

（1）上板旋转一周后，求球心在水平面内的位置。

（2）初始时球面最高点记为 \(A\)。上板转一周后，若矢量
\(\overrightarrow{CA}\) 恰好再次指向 \(+z\) 方向，求 \(r/l\)
的所有可能值。

\begin{figure}[H]
\centering
\includegraphics[width=0.821\linewidth,height=\textheight,keepaspectratio,alt={图2-1　旋转上板、固定下板与夹持球的三维几何关系}]{../figures/question-02.pdf}
\caption*{图2-1　旋转上板、固定下板与夹持球的三维几何关系}
\end{figure}

\subsection{一、题型定位}\label{ux4e00ux9898ux578bux5b9aux4f4d-1}

球在固定下板和转动上板之间纯滚。两个接触点的速度约束直接给出球心运动；追踪球面材料点时，进入随球心公转的坐标系可把转轴固定下来。

\subsection{二、知识点}\label{ux4e8cux77e5ux8bc6ux70b9-1}

\begin{enumerate}
\def\labelenumi{\arabic{enumi}.}
\tightlist
\item
  刚体上一点速度 \(\bm v_P=\bm v_C+\bm\omega\times\bm r_{P/C}\)。
\item
  两接触点的无滑动约束。
\item
  球心的半角速率公转。
\item
  旋转参考系中的角速度叠加 \(\bm\omega'=\bm\omega-\bm\Omega_f\)。
\item
  绕固定轴的有限转动与材料矢量复位条件。
\item
  整数绕转条件与离散几何参数。
\end{enumerate}

\subsection{三、物理图像}\label{ux4e09ux7269ux7406ux56feux50cf-1}

下接触点瞬时静止，上接触点随上板作绕 \(O\)
的圆周运动。上下接触点关于球心反对称，因此把两个无滑动方程相加时，球的自转项消失，球心速度恰好等于上接触点板速的一半。球心以角速度
\(\Omega/2\) 绕 \(O\) 公转。

球的水平角速度始终沿 \(OC\) 的反方向，因而在实验室系中不断转向；但在以
\(\Omega/2\) 绕 \(z\) 轴转动的随动系中，\(OC\)
固定，球的相对角速度也固定。材料点复位就转化为绕一根固定斜轴转过整数个整周。

\subsection{四、解题入口}\label{ux56dbux89e3ux9898ux5165ux53e3-1}

写出上、下接触点相对球心的位置

\[
\bm r_+=r\bm e_z,
\qquad
\bm r_-=-r\bm e_z.
\]

上板接触点速度为 \(\Omega\bm e_z\times\bm r_C\)，下板为零。两式相加求
\(\bm v_C\)，相减求 \(\bm\omega\)。随后选取以 \(\Omega/2\)
转动的坐标系，计算一周时间内相对转角。

\subsection{五、解析}\label{ux4e94ux89e3ux6790-1}

\begin{figure}[H]
\centering
\includegraphics[width=0.9\linewidth,height=\textheight,keepaspectratio,alt={图2-2　球心半速公转与随动系中的固定合成转轴}]{../figures/solution-02.pdf}
\caption*{图2-2　球心半速公转与随动系中的固定合成转轴}
\end{figure}

图2-2中，上下接触点确在同一竖直直径两端。左图记录球心半速公转，右图的固定斜轴属于随动系；材料点的姿态还要加回该坐标系绕
\(z\) 的转角。

\needspace{3\baselineskip}
\textbf{（一）两个接触点的无滑动方程}

设球心位置与球的角速度为

\[
\bm r_C=x\bm e_x+y\bm e_y, \tag{1}
\]

\[
\bm\omega=\omega_x\bm e_x+\omega_y\bm e_y,
\qquad \omega_z=0. \tag{2}
\]

球面上、下接触点的速度分别为

\[
\bm v_+=\bm v_C+\bm\omega\times(r\bm e_z), \tag{3}
\]

\[
\bm v_-=\bm v_C+\bm\omega\times(-r\bm e_z). \tag{4}
\]

下板固定且无滑动，所以

\[
\bm v_-=0. \tag{5}
\]

上板刚体转动，位于 \(\bm r_C\) 上方的接触点板速为

\[
\bm v_{\rm plate}=\Omega\bm e_z\times\bm r_C. \tag{6}
\]

上接触点无滑动给出

\[
\bm v_+=\Omega\bm e_z\times\bm r_C. \tag{7}
\]

把式（3）、（4）相加：

\[
2\bm v_C=\Omega\bm e_z\times\bm r_C. \tag{8}
\]

故

\[
\boxed{\bm v_C=\frac{\Omega}{2}\bm e_z\times\bm r_C}. \tag{9}
\]

该方程描述以角速度 \(\Omega/2\) 绕 \(O\) 的匀速圆周运动。

\needspace{3\baselineskip}
\textbf{（二）一周后的球心位置}

上板旋转一周所需时间

\[
T=\frac{2\pi}{\Omega}. \tag{10}
\]

在这段时间内，球心绕 \(O\) 转过

\[
\Delta\varphi_C=\frac\Omega2T=\pi. \tag{11}
\]

初始 \(\bm r_C(0)=l\bm e_x\)，因此

\[
\boxed{\bm r_C(T)=-l\bm e_x}, \tag{12}
\]

即球心坐标为

\[
\boxed{(-l,0)}. \tag{13}
\]

\needspace{3\baselineskip}
\textbf{（三）球的瞬时角速度}

由下接触点条件式（5）：

\[
\bm v_C-r\bm\omega\times\bm e_z=0. \tag{14}
\]

代入式（9）：

\[
r\bm\omega\times\bm e_z
=\frac\Omega2\bm e_z\times\bm r_C. \tag{15}
\]

展开叉积：

\[
\bm\omega\times\bm e_z
=\omega_y\bm e_x-\omega_x\bm e_y, \tag{16}
\]

\[
\bm e_z\times\bm r_C
=-y\bm e_x+x\bm e_y. \tag{17}
\]

逐分量比较：

\[
\omega_y=-\frac{\Omega}{2r}y,
\qquad
\omega_x=-\frac{\Omega}{2r}x. \tag{18}
\]

所以

\[
\boxed{\bm\omega=-\frac{\Omega}{2r}\bm r_C}. \tag{19}
\]

它是水平的，且始终指向 \(C\to O\)，与题设 \(\omega_z=0\) 一致。

\needspace{3\baselineskip}
\textbf{（四）随动坐标系中的固定转轴}

取坐标系 \(S'\) 以角速度

\[
\bm\Omega_f=\frac\Omega2\bm e_z \tag{20}
\]

绕实验室系转动。球心在 \(S'\) 中固定于
\(l\bm e_x'\)，因此式（19）也固定为

\[
\bm\omega=-\frac{\Omega l}{2r}\bm e_x'. \tag{21}
\]

刚体相对 \(S'\) 的角速度为

\[
\begin{aligned}
\bm\omega'
&=\bm\omega-\bm\Omega_f\\
&=-\frac\Omega2
\left(\frac lr\bm e_x'+\bm e_z\right).
\end{aligned} \tag{22}
\]

其大小为

\[
|\bm\omega'|
=\frac\Omega2\sqrt{1+\frac{l^2}{r^2}}. \tag{23}
\]

在一周时间内，球相对随动系转过

\[
\Phi=|\bm\omega'|T
=\pi\sqrt{1+\frac{l^2}{r^2}}. \tag{24}
\]

随动系自身绕 \(z\) 轴转过 \(\pi\)，但 \(z\) 方向不变。要让初始的
\(\overrightarrow{CA}=r\bm e_z\) 在末态仍指向 \(+z\)，只需它在 \(S'\)
中绕斜轴 \(\bm\omega'\) 后回到自身。

由于 \(l>0\)，\(\bm\omega'\) 不平行于
\(\bm e_z\)，材料矢量回到自身的条件只能是完整转过整数周：

\[
\Phi=2\pi n,
\qquad n=1,2,3,\ldots \tag{25}
\]

代入式（24）：

\[
\sqrt{1+\frac{l^2}{r^2}}=2n. \tag{26}
\]

整理得

\[
\boxed{\frac rl=\frac1{\sqrt{4n^2-1}}},
\qquad n=1,2,3,\ldots \tag{27}
\]

最前几个值为 \(1/\sqrt3,1/\sqrt{15},1/\sqrt{35},\ldots\)。

\needspace{3\baselineskip}
\textbf{有限转动复核}

式（22）给出固定单位轴

\[
\hat{\bm u}
=-\frac{(l/r)\bm e_x'+\bm e_z}{\sqrt{1+l^2/r^2}}. \tag{28}
\]

由 Rodrigues 公式，\(\bm e_z\) 绕 \(\hat{\bm u}\) 转角 \(\Phi\) 后为

\[
\bm e_z^{\,\rm f}
=\bm e_z\cos\Phi
+(\hat{\bm u}\times\bm e_z)\sin\Phi
+\hat{\bm u}(\hat{\bm u}\cdot\bm e_z)(1-\cos\Phi). \tag{29}
\]

因为 \(\hat{\bm u}\) 不平行于 \(\bm e_z\)，要使式（29）等于
\(\bm e_z\)，必须有 \(\sin\Phi=0\) 且 \(1-\cos\Phi=0\)，即
\(\Phi=2\pi n\)。这严格排除了``转过奇数个半周也可能复位''的误判。

\needspace{3\baselineskip}
\textbf{用旋转矩阵证明材料点必须整周回归}

随动系的固定转轴单位矢量为

\[
\hat{\bm n}={-(l/r)\hat{\bm e}_r-\hat{\bm e}_z\over\sqrt{1+l^2/r^2}}.
\]

初始顶部材料点的方向为 \(\hat{\bm e}_z\)。绕该轴转过 \(\Phi\)
后，Rodrigues 公式给出

\[
\bm u=\hat{\bm e}_z\cos\Phi
+(\hat{\bm n}\times\hat{\bm e}_z)\sin\Phi
+\hat{\bm n}(\hat{\bm n}\cdot\hat{\bm e}_z)(1-\cos\Phi).
\]

其竖直分量为

\[
u_z=\cos\Phi+n_z^2(1-\cos\Phi).
\]

要求回到顶部，必须 \(u_z=1\)。因 \(l>0\) 时 \(n_z^2<1\)，得到
\((1-n_z^2)(1-\cos\Phi)=0\)，即
\(\Phi=2\pi n\)。此时另外两个分量也自动为零。

上板一周后随动坐标系绕 \(z\) 转过 \(\pi\)，顶部方向不受这一额外绕 \(z\)
旋转影响，因此上述整周条件确能用于实验室的顶部回归。它并不声称任意初始材料点都满足同一回归条件；原题的``顶部点''使第二步旋转可以简化。

\subsection{六、关键洞察}\label{ux516dux5173ux952eux6d1eux5bdf-1}

\begin{enumerate}
\def\labelenumi{\arabic{enumi}.}
\tightlist
\item
  上下接触点约束相加消去自转，相减消去平动，是本题最快的代数组织方式。
\item
  球心公转角速度只有上板的一半，这是两接触面速度的算术平均。
\item
  实验室系中的水平自转轴在转动；选择与球心同速公转的坐标系后，它与坐标系转动角速度合成为固定斜轴。
\item
  第二问不是要求球的整体姿态完全回到初态，只要求一根特定材料矢量复位；但由于该矢量不平行于固定转轴，仍必须转过整周。
\end{enumerate}

\subsection{七、易错点}\label{ux4e03ux6613ux9519ux70b9-1}

\begin{enumerate}
\def\labelenumi{\arabic{enumi}.}
\tightlist
\item
  把上接触点的板速写成
  \(\Omega\bm e_z\times(r\bm e_z)\)；应使用接触点相对转轴 \(O\)
  的水平位置 \(\bm r_C\)。
\item
  认为球心也以角速度 \(\Omega\) 公转，漏掉因子 \(1/2\)。
\item
  由式（19）直接把 \(|\bm\omega|T\)
  当作总转角；转轴方向在实验室系中随时间变化。
\item
  在旋转系中只保留水平角速度，漏掉 \(-\Omega\bm e_z/2\) 的坐标系转动项。
\item
  把复位条件写成 \(\Phi=n\pi\)，没有检查奇数倍 \(\pi\) 会把 \(\bm e_z\)
  转到另一方向。
\end{enumerate}

\subsection{八、课堂讲解建议}\label{ux516bux8bfeux5802ux8bb2ux89e3ux5efaux8bae-1}

先给出``一只球夹在两条不同速度传送带之间''的二维类比，让学生猜球心速度是两带速度的平均值。再回到本题，用矢量式证明。第一问完成后，要明确区分球心的公转角速度与球自身的角速度。

第二问在黑板上保留两个坐标系：实验室系画出转动的径向轴，随动系画出固定的
\(\bm e_x'\) 和 \(\bm e_z\)。把 \(\bm\omega'\)
的水平、竖直两个分量画成直角三角形，再求其模。

可让学生取 \(n=1\)，算得 \(r/l=1/\sqrt3\)，再代入式（24）得到
\(\Phi=2\pi\)；取一个不满足离散条件的数值，则材料点不能复位。

\subsection{九、变式训练}\label{ux4e5dux53d8ux5f0fux8badux7ec3-1}

\begin{enumerate}
\def\labelenumi{\arabic{enumi}.}
\tightlist
\item
  若下板也以角速度 \(\Omega_-\) 转动，上板角速度为
  \(\Omega_+\)，求球心公转角速度。
\item
  去掉``\(\omega_z=0\)''条件，设球初始另有竖直自转
  \(\omega_{z0}\)，重新给出复位条件。
\item
  上板只转过角度 \(\Theta\)，求球心位置并写出材料矢量复位方程。
\item
  用旋转矩阵数值计算 \(n=1\) 时 \(\bm e_z\) 的末态。
\end{enumerate}

\needspace{3\baselineskip}
\textbf{变式解答：第一、第二个材料点回归条件}

球心半速公转，上板一周 \(2\pi/\Omega\)
后球心位于对径点。随动系的转轴固定，角速度大小为

\[
|\bm\omega'|={\Omega\over2}\sqrt{1+{l^2\over r^2}}.
\]

该时间内随动系恰好绕 \(z\) 轴转过 \(\pi\)。初、末顶部方向均沿 \(z\)，绕
\(z\) 的这一步不会改变目标方向，故材料点回到顶部要求随动旋转整周：

\[
\pi\sqrt{1+{l^2\over r^2}}=2\pi n.
\]

由此

\[
{r\over l}={1\over\sqrt{4n^2-1}}.
\]

第一、第二个允许半径分别为
\(l/\sqrt3\)、\(l/\sqrt{15}\)。一般点若恰在随动转轴上，旋转并不改变其方向，但原题标记在顶部，转轴含非零径向分量，因此不属于这种例外。

检查滚动约束时，先由下接触点
\(\bm v_C+\bm\omega\times(-r\hat{\bm z})=0\) 和上接触点速度求出
\(\bm\omega\)。直接把材料点当成始终绕竖直轴转动，会得到球心位置正确、姿态却错误的结果。

\newpage

\section{第 3 题}\label{ux7b2c-3-ux9898}

如图为无限角点相接三角电阻网络，每条小三角形边的电阻均为 \(r\)。节点
\(A,B,C\)
是同一个小三角形的三个顶点，由三重旋转对称互相对应。原图是四配位的角共享三角网络：每个节点属于两个三角形，周围另有六边形空隙；它不是六配位的完整三角晶格。

（1）求相邻节点 \(A,B\) 间的等效电阻 \(r_{AB}\)。

（2）在 \(A,B,C\) 接入三相交流电源

\[
u_A=u_0\cos\omega t,\quad
u_B=u_0\cos\left(\omega t+\frac{2\pi}{3}\right),\quad
u_C=u_0\cos\left(\omega t+\frac{4\pi}{3}\right),
\]

无穷远接地。求从 \(A\) 端流入网络的电流 \(I_A(t)\)。

（3）把无限网络视为 \(ABC\) 三端网络并等效成星形网络，求 \(Z_A\)。

（4）若直接连接 \(A,B\) 的那一条电阻断开，求新的星形等效阻抗
\(Z_A',Z_C'\)。

\begin{figure}[H]
\centering
\includegraphics[width=0.9\linewidth,height=\textheight,keepaspectratio,alt={图3-1　角点相接无限三角网络与三端星形等效表示}]{../figures/question-03.pdf}
\caption*{图3-1　角点相接无限三角网络与三端星形等效表示}
\end{figure}

\subsection{一、题型定位}\label{ux4e00ux9898ux578bux5b9aux4f4d-2}

原图由角点相接的三角形组成，每个节点接四条边。最近邻电阻可由周期网络的对称性求出，三端等效星形再处理三相电压和断边后的网络。

\subsection{二、知识点}\label{ux4e8cux77e5ux8bc6ux70b9-2}

\begin{enumerate}
\def\labelenumi{\arabic{enumi}.}
\tightlist
\item
  无限均匀网络中从一点注入、流向无穷远的等分流。
\item
  线性电阻网络的叠加原理。
\item
  相量表示与三相对称关系 \(\tilde u_A+\tilde u_B+\tilde u_C=0\)。
\item
  对称三端网络的星形等效。
\item
  串并联反演与断边后的不对称星形参数。
\item
  电阻网络结果的有限截断数值核验。
\end{enumerate}

\subsection{三、物理图像}\label{ux4e09ux7269ux7406ux56feux50cf-2}

图3-1重画了两个三角形在角点共享、六边形留空的实际拓扑。每个节点有四条出边，在无限网络对称性下对单源问题等价。从
\(A\) 向无穷远注入电流 \(I\) 时，离开 \(A\) 的四条边各流
\(I/4\)；从无穷远汇入 \(B\) 时，进入 \(B\) 的四条边也各流
\(I/4\)。将两种场叠加为``\(A\) 注入、\(B\) 流出''，直接边 \(AB\)
上的电流同向相加为 \(I/2\)，从而立即读出 \(r_{AB}\)。

三相激励时，三个端电压相量和为零，星形中点电势也为零。每一相只通过自己的星形支路流向公共中点，问题退化为单个电阻上的欧姆定律。

\subsection{四、解题入口}\label{ux56dbux89e3ux9898ux5165ux53e3-2}

第一问分成两个辅助状态：状态甲为 \(A\) 注入
\(I\)、无穷远吸收；状态乙为无穷远注入 \(I\)、\(B\)
吸收。两个状态叠加后，除 \(A,B\) 外各节点仍满足基尔霍夫电流定律。

第二、三问直接使用由第一问得到的端间电阻。第四问把``断边后的剩余网络''画成未知星形，再把阻值为
\(r\) 的 \(AB\) 直连边并联回去。

\subsection{五、解析}\label{ux4e94ux89e3ux6790-2}

\begin{figure}[H]
\centering
\includegraphics[width=0.9\linewidth,height=\textheight,keepaspectratio,alt={图3-2　断去 AB 直连边后的星形参数与三相电流关系}]{../figures/solution-03.pdf}
\caption*{图3-2　断去 \(AB\) 直连边后的星形参数与三相电流关系}
\end{figure}

图3-2中，题面是真实的四配位网络。解析左图只画断去 \(AB\)
后的等效星形，右图画完整网络的三相驱动；两种电路状态各用自己的星形支路值。

\needspace{3\baselineskip}
\textbf{（一）相邻节点的等效电阻}

在状态甲中，从 \(A\) 注入总电流 \(I\)，电流最终流向无穷远。由于 \(A\)
的四条出边在网络对称性下等价，每条边电流为

\[
I_{AB}^{(1)}=\frac I4. \tag{1}
\]

在状态乙中，从无穷远向 \(B\) 汇入总电流 \(I\)。进入 \(B\)
的四条边等价，直接边 \(AB\) 上的电流同样为

\[
I_{AB}^{(2)}=\frac I4, \tag{2}
\]

方向仍是 \(A\to B\)。叠加两状态后，直接边上的电流为

\[
I_{AB}=I_{AB}^{(1)}+I_{AB}^{(2)}=\frac I2. \tag{3}
\]

该边电压降

\[
U_{AB}=rI_{AB}=\frac{Ir}{2}. \tag{4}
\]

所以

\[
\boxed{r_{AB}=\frac{U_{AB}}I=\frac r2}. \tag{5}
\]

这里的``每点四条等价出边''由角点相接网络的实际连接关系决定；若误看成每个几何交叉处都电连接的普通三角晶格，配位数会变成
\(6\)，答案也会改变。

\needspace{3\baselineskip}
\textbf{（二）三相电源下的端电流与星形参数}

用复数相量表示三相电压：

\[
\tilde u_A=u_0,\quad
\tilde u_B=u_0e^{i2\pi/3},\quad
\tilde u_C=u_0e^{i4\pi/3}. \tag{6}
\]

它们满足

\[
\tilde u_A+\tilde u_B+\tilde u_C=0. \tag{7}
\]

由三重旋转对称性，星形三支路相等：

\[
Z_A=Z_B=Z_C=Z. \tag{8}
\]

当只在 \(A,B\) 间测量时，\(C\) 端开路，故

\[
r_{AB}=Z_A+Z_B=2Z. \tag{9}
\]

代入式（5）：

\[
\boxed{Z_A=Z_B=Z_C=\frac r4}. \tag{10}
\]

星形中点相量为三个端电压的平均值：

\[
\tilde u_N
=\frac{\tilde u_A+\tilde u_B+\tilde u_C}{3}=0. \tag{11}
\]

于是

\[
\tilde I_A
=\frac{\tilde u_A-\tilde u_N}{r/4}
=\frac{4u_0}{r}. \tag{12}
\]

恢复时间形式：

\[
\boxed{I_A(t)=\frac{4u_0}{r}\cos\omega t}. \tag{13}
\]

由于网络只有纯电阻，电流与本相电压同相。

\needspace{3\baselineskip}
\textbf{（三）断开 \(AB\) 边后的星形参数}

记断边后剩余网络的星形参数为

\[
Z_A'=Z_B'=a,\qquad Z_C'=c. \tag{14}
\]

其中 \(A,B\) 对称性仍保留。若把阻值 \(r\) 的直接边重新接回，原网络在
\(A,B\) 间的等效电阻为

\[
\frac r2=r\parallel(2a)
=\frac{2ar}{r+2a}. \tag{15}
\]

交叉相乘：

\[
\frac r2(r+2a)=2ar, \tag{16}
\]

\[
r+2a=4a. \tag{17}
\]

所以

\[
\boxed{a=Z_A'=Z_B'=\frac r2}. \tag{18}
\]

再看原网络的 \(A,C\) 两端。接回 \(AB\) 边后，从星形中心到 \(A\)
有两条路线：直接经过 \(a\)，或先经过 \(a\) 到 \(B\) 再经过电阻 \(r\) 到
\(A\)。因此中心到 \(A\) 的等效为

\[
a\parallel(a+r)
=\frac{a(a+r)}{2a+r}. \tag{19}
\]

从 \(C\) 到中心串联 \(c\)，故

\[
\frac r2
=c+\frac{a(a+r)}{2a+r}. \tag{20}
\]

代入 \(a=r/2\)：

\[
\frac{a(a+r)}{2a+r}
=\frac{(r/2)(3r/2)}{2r}
=\frac{3r}{8}. \tag{21}
\]

于是

\[
c=\frac r2-\frac{3r}{8}=\frac r8. \tag{22}
\]

最终

\[
\boxed{Z_A'=Z_B'=\frac r2,\qquad Z_C'=\frac r8}. \tag{23}
\]

\needspace{3\baselineskip}
\textbf{功率复核}

三相平衡时每一星形支路电阻为 \(r/4\)，每相电压有效值为
\(u_0/\sqrt2\)，总平均功率

\[
\bar P
=3\frac{(u_0/\sqrt2)^2}{r/4}
=\frac{6u_0^2}{r}. \tag{24}
\]

另一方面每相电流峰值为 \(4u_0/r\)，每相平均功率为
\(2u_0^2/r\)，三相之和相同，说明电流幅值与相位自洽。

\needspace{3\baselineskip}
\textbf{周期逼近与三端中性模}

二维无限网向无穷远送入净电流时，绝对电势有对数发散，不能把有限的 \(V_A\)
当作已有结论。上面的辅助单源应先在有限网络或周期网络上定义，再取电势差的极限；源汇叠加后的两端问题净注入为零，电势差具有有限极限。

取由 \(L\times L\) 个单元构成、每单元三个节点的周期角共享三角网络，共有
\(N=3L^2\) 个节点、\(2N\) 条等价边。Foster 关系给出

\[
{r_{AB}(L)\over r}={N-1\over2N}\longrightarrow{1\over2}.
\]

平衡三相满足总电流和总电压相量的中性条件，公共模不被激励。由三端对称性，两端电阻等于两条星形支路之和，得到
\(Z=r/4\)；这一步不需要为发散的单源绝对电势赋值。实际图中的四条边计数与周期结果相互核验。

\needspace{3\baselineskip}
\textbf{变式解答：三端导纳矩阵与功率}

完整网络等效星形的三支路均为 \(r/4\)。悬空中心点电势是三端平均：

\[
U_O={U_A+U_B+U_C\over3}.
\]

因此端口电流为

\[
I_A={4\over3r}(2U_A-U_B-U_C),
\]

其他两式由循环换标号得到。三端电流之和为零，给出导纳矩阵的行和检验。三相驱动满足
\(U_A+U_B+U_C=0\)，所以

\[
I_A={4U_A\over r},\qquad
P(t)={4\over r}(U_A^2+U_B^2+U_C^2).
\]

三相等幅正弦的平方和为 \(3u_0^2/2\)，总耗散功率恒为
\(6u_0^2/r\)。这是一项整体检验，不依赖如何画局部电阻。

删去 \(AB\) 后，星形支路变为
\(r/2,r/2,r/8\)；在两端口实验中第三端悬空，等效电阻为
\(R_{AB}'=r\)、\(R_{AC}'=5r/8\)。把第三端悬空的串联结果用于三端同时加电是不合适的，必须先求中心点电势再计算电流。

\subsection{六、关键洞察}\label{ux516dux5173ux952eux6d1eux5bdf-2}

\begin{enumerate}
\def\labelenumi{\arabic{enumi}.}
\tightlist
\item
  无限网络的难点被局部对称性消去；只要知道源点的等价出边数，就能得到相邻边电流。
\item
  叠加的两个辅助状态必须在除源、汇点外的所有节点满足同一基尔霍夫条件。
\item
  对称三相电压相量和为零，使星形中点成为虚地；这不是因为中点真实接地。
\item
  断边后虽然整个网络不再三重对称，但 \(A\leftrightarrow B\)
  镜面对称仍在，因此只有两个独立参数。
\end{enumerate}

\subsection{七、易错点}\label{ux4e03ux6613ux9519ux70b9-2}

\begin{enumerate}
\def\labelenumi{\arabic{enumi}.}
\tightlist
\item
  把图中几何交叉都当成电连接节点，错误地采用六配位三角晶格。
\item
  状态甲、乙的 \(AB\) 电流方向判断相反，导致相减而不是相加。
\item
  把 \(r_{AB}=r/2\) 直接当作星形单支路，漏掉两支路串联关系。
\item
  三相相量的指数符号写错，破坏相量和为零。
\item
  第四问仍令三支路相等，忽略断边后 \(C\) 与 \(A,B\) 的不等价。
\end{enumerate}

\subsection{八、课堂讲解建议}\label{ux516bux8bfeux5802ux8bb2ux89e3ux5efaux8bae-2}

先画源点附近的四条出边，让学生口头完成``每边
\(I/4\)''。然后把``\(A\to\infty\)''和``\(\infty\to B\)''用两种颜色叠在同一张图上，直接连接
\(A,B\)
的边上，两股电流同向相加。无穷远只是汇流边界，不需要知道每一条远处支路的电流。

三相部分先验证
\(1+e^{i2\pi/3}+e^{i4\pi/3}=0\)，再画星形中点。第四问先只求 \(A,B\)
间，把未知量降到 \(a\)；得到 \(a=r/2\) 后再处理 \(A,C\)。

若条件允许，可用有限截断网络做数值检验。随着截断尺度增大，\(r_{AB}\)
应向 \(r/2\) 收敛。数值验证只能检查结论，不能替代对称证明。

\subsection{九、变式训练}\label{ux4e5dux53d8ux5f0fux8badux7ec3-2}

\begin{enumerate}
\def\labelenumi{\arabic{enumi}.}
\tightlist
\item
  若网络每个节点有 \(z\) 条等价出边，证明相邻节点等效电阻为 \(2r/z\)。
\item
  三相电压幅值不等时，求星形中点电势与各相电流。
\item
  若断开的是 \(A,C\) 边，利用置换对称性写出新的星形参数。
\item
  在 \(AB\) 直连边上串入电阻 \(R\)，求修改后三端网络参数。
\end{enumerate}

\needspace{3\baselineskip}
\textbf{不用``无穷远电势''为四配位网提供证明}

选一个包含 \(N=3L^2\) 个节点的周期 Kagome
网络，保持三角形和六角孔的构造。每个节点有四条边，无向边共有 \(2N\)
条。各边由晶格平移与旋转对称相互映射，Foster 求和给出

\[
2N{R_{\rm edge}(N)\over r}=N-1.
\]

因此

\[
R_{\rm edge}(N)={r\over2}\left(1-{1\over N}\right)
\longrightarrow{r\over2}.
\]

用单位源汇分别接在相邻两个点，解离散方程

\[
\sum_{j\sim i}(V_i-V_j)=I_i r
\]

即可计算有限尺寸值。\(L=12\) 时 \(N=432\)，得到
\(R/r=431/864\simeq0.4988426\)；它不是不精确数值误差，而是已知的周期有限尺寸修正。

随后用三端对称星形的 \(2Z=R_{\rm edge}\)，无限尺寸得到
\(Z=r/4\)。这个过程也检查原图确为四配位；若误建成六配位三角晶格，边数和最近邻电阻都会改变。周期网络只用于检验源汇电压差，不声称二维单源相对无穷远的绝对电势有限。

\needspace{3\baselineskip}
\textbf{删去直连边后的导纳核验}

完整三端星形在去掉共同电势自由度后，导纳为

\[
Y={4\over r}\left(I-{\bm1\bm1^{\rm T}\over3}\right).
\]

删去 \(AB\) 电阻，相当于减去该边的导纳矩阵。取
\(\bm b=(1,-1,0)^{\rm T}\)，则

\[
Y'=Y-{\bm b\bm b^{\rm T}\over r}.
\]

在 \(\bm b\) 方向，原特征值 \(4/r\) 减去 \(2/r\)，余下
\(2/r\)；与其正交的零和方向仍为 \(4/r\)。因此

\[
R'_{AB}={|\bm b|^2\over2/r}=r.
\]

对 \(AC\) 源汇向量 \(\bm a=(1,0,-1)\)，在 \(\bm b\) 方向的投影平方为
\(1/2\)，正交部分平方为 \(3/2\)，所以

\[
R'_{AC}={{1/2}\over2/r}+{{3/2}\over4/r}={5r\over8}.
\]

由 \(Z_A+Z_B=r\)、\(Z_A+Z_C=Z_B+Z_C=5r/8\) 解得
\(Z_A=Z_B=r/2\)、\(Z_C=r/8\)。该计算使用同一三端网络的导纳变化，避免凭删边图的外观猜测各支路如何改变；结果与图3-2及并联求
\(R'_{AB}\) 的路线一致。

\newpage

\section{第 4 题}\label{ux7b2c-4-ux9898}

如图（a），质量为 \(m\)、半径为 \(R\) 的均匀绝缘圆环铰接于环上一点
\(A\)，可在圆环所在的水平面内自由转动。质量为 \(m\)、电荷量为 \(q\)
的小珠套在圆环上，与圆环间无摩擦。空间中存在竖直向下的匀强磁场
\(\bm B\)。

（1）若小珠相对圆环静止，并随圆环绕 \(A\) 做圆周运动，求共同角速度
\(\omega_1\)。小珠到 \(A\) 的距离记为 \(r_1\)，其中 \(0<r_1<2R\)。

（2）使小珠起初位于 \(A\)，圆环以角速度 \(\omega_0\) 转动，且
\(\omega_0\)
大于第（1）问所得角速度。给小珠一个微小扰动后，小珠沿环滑动。若小珠能够第一次再次与圆环相对静止，求此时小珠到
\(A\) 的距离 \(r\) 和圆环角速度 \(\omega_2\)。

（3）如图（b），把小珠固定在质量同为 \(m\)、半径为 \(R/2\)
的均匀小绝缘环上；小环在大环内无阻力地纯滚动。初始时两环在 \(A\)
点相切并相对静止，小珠也位于 \(A\)。大环初角速度为

\[
\omega_0'=0.53\frac{qB}{m}.
\]

给小环微小扰动，使其圆心相对大环圆心顺时针运动。若两环能够第一次再次相对静止，求小环圆心转过的角度
\(\alpha\) 以及此时两环的共同角速度 \(\omega_2\)。

\begin{figure}[H]
\centering
\includegraphics[width=0.9\linewidth,height=\textheight,keepaspectratio,alt={图4-1　铰接圆环上的带电小珠及内滚小环的几何关系}]{../figures/question-04.pdf}
\caption*{图4-1　铰接圆环上的带电小珠及内滚小环的几何关系}
\end{figure}

\subsection{一、题型定位}\label{ux4e00ux9898ux578bux5b9aux4f4d-3}

磁场虽不做功，却会改变带电小珠绕铰点的机械角动量。守恒的是含磁矢势项的正则角动量；小环在大环内滚动时，还须先求准小珠相对大环的位置。

\subsection{二、知识点}\label{ux4e8cux77e5ux8bc6ux70b9-3}

\begin{enumerate}
\def\labelenumi{\arabic{enumi}.}
\tightlist
\item
  洛伦兹力提供向心力的条件。
\item
  磁场不做功与机械能守恒。
\item
  关于固定点的磁力矩及其时间积分。
\item
  均匀细圆环关于环上一点的转动惯量。
\item
  内滚动约束与圆心转角、环自转角的关系。
\item
  瞬时相对静止时系统总动能的统一表示。
\item
  几何可达条件与物理解的筛选。
\end{enumerate}

\subsection{三、物理图像}\label{ux4e09ux7269ux7406ux56feux50cf-3}

小珠运动时，洛伦兹力方向总在水平面内且垂直于小珠速度。它不改变速率，却能改变系统关于
\(A\) 的机械角动量。把小珠位置矢量写成 \(\bm r\)，磁力矩的 \(z\)
分量恰好是
\(\tfrac12qB\,\mathrm d(r^2)/\mathrm dt\)；因此机械角动量的变化可被一个只依赖位置的项吸收。

第三问中，小环半径是大环的一半。内切时两圆心间距为
\(R/2\)，纯滚动使小环相对大环的自转角与圆心公转角大小相等、方向相反；小珠与当前接触半径的夹角则为
\(2\alpha\)。由于带电小珠固定在小环上并从内切点 \(A\)
出发，它在整个运动中都落在直线 \(AO\) 上，到 \(A\) 的距离化为
\(R(1-\cos\alpha)\)。

\subsection{四、解题入口}\label{ux56dbux89e3ux9898ux5165ux53e3-3}

先从磁力矩直接推导广义角动量守恒式，而不要把``磁场不做功''误写成``磁场没有力矩''。第（2）问在初、末两个相对静止时刻使用广义角动量守恒和机械能守恒，消去
\(\omega_2\) 即可求 \(r\)。

第（3）问的关键不是重新列复杂的瞬时速度，而是只比较两个``相对静止''的时刻。此时整套构形瞬时以共同角速度
\(\omega\) 绕 \(A\) 转动，故总动能可写为
\(\tfrac12I_{\rm eff}(\alpha)\omega^2\)。把大环、小环和小珠对 \(A\)
的等效转动惯量逐项相加，即能复用同一组守恒方程。

\subsection{五、解析}\label{ux4e94ux89e3ux6790-3}

\begin{figure}[H]
\centering
\includegraphics[width=0.9\linewidth,height=\textheight,keepaspectratio,alt={图4-2　广义角动量守恒与内滚构形中的关键几何量}]{../figures/solution-04.pdf}
\caption*{图4-2　广义角动量守恒与内滚构形中的关键几何量}
\end{figure}

图4-2中，左图用点号标出垂直纸面的机械角动量，右图按相切条件把小环完整放在大环内。小珠在
\(AO\) 上的距离 \(R(1-\cos\alpha)\) 随后进入正则角动量与动能方程。

\needspace{3\baselineskip}
\textbf{（一）广义角动量守恒式的来源}

取铰点 \(A\) 为坐标原点，\(z\)
轴正方向与所选角速度正方向一致。带电小珠的位置与速度分别为

\[
\bm r=(x,y,0),\qquad \bm v=(\dot x,\dot y,0).
\]

按题设磁场方向和正方向约定，洛伦兹力在水平面的分量可写成

\[
\bm F_B=q\bm v\times\bm B.
\]

磁力关于 \(A\) 的力矩为 \(\bm\tau_B=\bm r\times\bm F_B\)。其 \(z\)
分量为

\[
\begin{aligned}
\tau_{B,z}
&=xF_y-yF_x\\
&=qB(x\dot x+y\dot y)\\
&=\frac{qB}{2}\frac{\mathrm d}{\mathrm dt}(x^2+y^2)\\
&=\frac{qB}{2}\frac{\mathrm d r^2}{\mathrm dt}. 
\end{aligned}
\tag{1}
\]

系统关于固定点 \(A\) 的机械角动量满足

\[
\frac{\mathrm dL_z}{\mathrm dt}=\tau_{B,z}. \tag{2}
\]

把式（1）代入式（2）并移项：

\[
\frac{\mathrm d}{\mathrm dt}\left(L_z-\frac12qBr^2\right)=0. \tag{3}
\]

所以

\[
\boxed{L_z-\frac12qBr^2=\text{常量}}. \tag{4}
\]

铰链作用力通过 \(A\)，关于 \(A\)
的力矩为零；圆环与小珠之间的约束力是系统内力，成对计入整套系统后不改变式（4）。另一方面

\[
\bm F_B\cdot\bm v=q(\bm v\times\bm B)\cdot\bm v=0, \tag{5}
\]

铰链点不移动，理想约束也不做净功，因而系统机械动能守恒。

\needspace{3\baselineskip}
\textbf{（二）第（1）问：共同匀速转动}

若小珠与圆环无相对运动，则小珠绕 \(A\) 做半径为 \(r_1\)
的匀速圆周运动。洛伦兹力的大小为

\[
F_B=qvB=q(\omega_1r_1)B. \tag{6}
\]

题设所需转向下，洛伦兹力指向 \(A\)，充当向心力：

\[
q\omega_1r_1B=m\omega_1^2r_1. \tag{7}
\]

由于 \(r_1>0\) 且讨论非零转动，可约去 \(\omega_1r_1\)：

\[
\boxed{\omega_1=\frac{qB}{m}}. \tag{8}
\]

式（8）与 \(r_1\) 无关，因此几何允许的任意

\[
\boxed{0<r_1<2R} \tag{9}
\]

都可满足。若 \(q\)
为负或反向选取转动正方向，应把式（8）理解为带符号的回旋角速度；其大小为
\(|q|B/m\)。

\needspace{3\baselineskip}
\textbf{（三）第（2）问：第一次再次相对静止}

均匀细圆环绕过环上一点 \(A\)、垂直环面的轴的转动惯量由平行轴定理得到

\[
I_A=I_O+mR^2=mR^2+mR^2=2mR^2. \tag{10}
\]

初态小珠在 \(A\)，故 \(r=0\)，小珠对 \(A\)
的角动量和动能均为零。初态机械角动量与动能分别为

\[
L_i=2mR^2\omega_0,\qquad
E_i=\frac12(2mR^2)\omega_0^2. \tag{11}
\]

末态小珠在距 \(A\) 为 \(r\) 的位置，并与圆环瞬时共同以 \(\omega_2\)
转动。此时对 \(A\) 的总转动惯量为

\[
I_f=2mR^2+mr^2. \tag{12}
\]

由式（4）得

\[
2mR^2\omega_0
=(2mR^2+mr^2)\omega_2-\frac12qBr^2. \tag{13}
\]

由机械能守恒得

\[
2mR^2\omega_0^2
=(2mR^2+mr^2)\omega_2^2. \tag{14}
\]

为完整展示消元过程，令

\[
X=2R^2+r^2. \tag{15}
\]

约去式（14）中的 \(m\) 后有

\[
\omega_2=\omega_0\sqrt{\frac{2R^2}{X}}. \tag{16}
\]

取与初态同向的正根，因为第一次相对静止前系统尚未整体反转。把式（16）代回式（13），再结合
\(r^2=X-2R^2\)，化简可得

\[
\sqrt{2R^2X}\,m\omega_0
=2mR^2\omega_0+\frac12qB(X-2R^2). \tag{17}
\]

也可先用式（14）把 \(I_f\omega_2\) 写成
\(I_i\omega_0^2/\omega_2\)。这样式（13）变为

\[
2mR^2\omega_0
=\frac{2mR^2\omega_0^2}{\omega_2}
-\frac12qBr^2. \tag{18}
\]

再由式（14）写出

\[
r^2=2R^2\left(\frac{\omega_0^2}{\omega_2^2}-1\right). \tag{19}
\]

代入式（18），约去公共因子并整理为

\[
(2m\omega_0-qB)\omega_2^2
-2m\omega_0^2\omega_2+qB\omega_0^2=0. \tag{20}
\]

左端可因式分解为

\[
(\omega_2-\omega_0)
\left[(2m\omega_0-qB)\omega_2-qB\omega_0\right]=0. \tag{20a}
\]

根 \(\omega_2=\omega_0\) 代回式（19）只给出 \(r=0\)，对应尚未离开 \(A\)
的初态，不是扰动后的第一次再次相对静止，故应舍去。非平凡根满足

\[
(2m\omega_0-qB)\omega_2=qB\omega_0. \tag{20b}
\]

因此

\[
\boxed{\omega_2=\frac{qB\omega_0}{2m\omega_0-qB}}. \tag{21}
\]

将式（21）代入式（19）：

\[
\begin{aligned}
r^2
&=2R^2\left[\frac{(2m\omega_0-qB)^2}{q^2B^2}-1\right]\\
&=2R^2\frac{4m^2\omega_0^2-4mqB\omega_0}{q^2B^2}\\
&=\frac{8mR^2\omega_0(m\omega_0-qB)}{q^2B^2}. 
\end{aligned}
\tag{22}
\]

于是

\[
\boxed{r=R\frac{\sqrt{8m\omega_0(m\omega_0-qB)}}{qB}}. \tag{23}
\]

根式要求 \(m\omega_0>qB\)，正好对应题设
\(\omega_0>\omega_1\)。还必须满足小珠在圆环上的几何条件 \(r\leq2R\)。令
\(a=m\omega_0/(qB)\)，则

\[
8a(a-1)\leq4
\quad\Longrightarrow\quad
1<a\leq\frac{1+\sqrt3}{2}. \tag{24}
\]

超过这一上界时，守恒方程给出的再次相对静止位置落到圆环几何范围之外，不能作为本问物理解。

\needspace{3\baselineskip}
\textbf{（四）第（3）问：内滚小环的几何关系}

设小环圆心为 \(O'\)，大环圆心为 \(O\)。小环半径为 \(R/2\)，内切约束给出

\[
OO'=R-\frac R2=\frac R2. \tag{25}
\]

设 \(OO'\) 相对初始方向顺时针转过 \(\alpha\)。取大环固连坐标，令
\(A=(0,0)\)、\(O=(R,0)\)。圆心从初始的左向半径顺时针转过 \(\alpha\) 后

\[
O'=\left(R-{R\over2}\cos\alpha,{R\over2}\sin\alpha\right).
\]

接触点相对大环的速度为零。对小环有
\(\bm v_{O'}+\bm\omega_{\rm rel}\times\bm r_{P/O'}=0\)，推出小环相对大环的自转角为
\(+\alpha\)，与圆心的 \(-\alpha\) 转角方向相反。故初始由 \(O'\) 指向
\(A\) 的材料半径变为

\[
\bm r_{\rm bead/O'}=-{R\over2}(\cos\alpha,\sin\alpha).
\]

把两个位置相加，得到小珠位置

\[
\boxed{\bm r_{\rm bead/A}=(R(1-\cos\alpha),0)}.
\]

因此小珠始终位于 \(AO\) 直线上，\(r=R(1-\cos\alpha)\)。有些推导写
\(R\alpha=(R/2)\beta\)，其中 \(\beta=2\alpha\)
指的是小珠与当前接触半径的相对角，并非小环相对大环的自转角。图4-2明确区分这两种角度。

记

\[
c=\cos\alpha. \tag{29}
\]

则 \(r=R(1-c)\)。

\needspace{3\baselineskip}
\textbf{（五）两个相对静止时刻的等效转动惯量}

初态对应 \(\alpha=0\)、\(c=1\)。大环关于 \(A\) 的转动惯量仍为
\(2mR^2\)。小环圆心初始距 \(A\) 为 \(R/2\)，小环自身绕心转动惯量为

\[
I_{O'}=m\left(\frac R2\right)^2=\frac14mR^2. \tag{30}
\]

故小环对 \(A\) 的转动惯量为

\[
I_{A,\text{小环},i}=I_{O'}+m\left(\frac R2\right)^2
=\frac12mR^2. \tag{31}
\]

小珠初始位于 \(A\)，贡献为零，所以

\[
I_i=2mR^2+\frac12mR^2=\frac52mR^2. \tag{32}
\]

末态两环瞬时再次相对静止，因此整个构形瞬时以共同角速度 \(\omega_2\) 绕
\(A\) 转动。由余弦定理

\[
\begin{aligned}
AO'^2
&=AO^2+OO'^2-2(AO)(OO')\cos\alpha\\
&=R^2+\frac{R^2}{4}-R^2c\\
&=R^2\left(\frac54-c\right). 
\end{aligned}
\tag{33}
\]

所以小环关于 \(A\) 的转动惯量为

\[
I_{A,\text{小环},f}
=\frac14mR^2+mR^2\left(\frac54-c\right)
=mR^2\left(\frac32-c\right). \tag{34}
\]

小珠的贡献为

\[
I_{A,\text{珠},f}=mr^2=mR^2(1-c)^2. \tag{35}
\]

加上大环：

\[
\begin{aligned}
I_f
&=2mR^2+mR^2\left(\frac32-c\right)+mR^2(1-c)^2\\
&=mR^2\left(c^2-3c+\frac92\right). 
\end{aligned}
\tag{36}
\]

\needspace{3\baselineskip}
\textbf{（六）守恒方程与数值求解}

在初、末两个相对静止时刻使用式（4）和机械能守恒：

\[
I_i\omega_0'
=I_f\omega_2-\frac12qBR^2(1-c)^2, \tag{37}
\]

\[
I_i\omega_0'^2=I_f\omega_2^2. \tag{38}
\]

由式（38）取同向根：

\[
\omega_2=\omega_0'\sqrt{\frac{I_i}{I_f}}. \tag{39}
\]

把式（39）代入式（37），并除以 \(I_i\omega_0'\)：

\[
1=\sqrt{\frac{I_f}{I_i}}-
\frac{qB}{5m\omega_0'}(1-c)^2. \tag{40}
\]

移项后平方，得到只含 \(c\) 的方程

\[
\boxed{
\left[1+\frac{qB}{5m\omega_0'}(1-c)^2\right]^2
=\frac25\left(c^2-3c+\frac92\right)}. \tag{41}
\]

代入 \(\omega_0'=0.53qB/m\)：

\[
\left[1+\frac{(1-c)^2}{2.65}\right]^2
=\frac25\left(c^2-3c+\frac92\right). \tag{42}
\]

在第一次离开初态后的物理解区间 \(-1<c<1\)
中数值求根，并剔除由平方引入的不满足式（40）的根，得到

\[
c=0.1335082176. \tag{43}
\]

因此

\[
\boxed{\alpha=\arccos c=82.3276^\circ}. \tag{44}
\]

再由式（39）：

\[
\begin{aligned}
\omega_2
&=0.53\frac{qB}{m}
\sqrt{\frac{5/2}{c^2-3c+9/2}}\\
&=0.4129901\frac{qB}{m}. 
\end{aligned}
\tag{45}
\]

故

\[
\boxed{\omega_2\approx0.4130\frac{qB}{m}}. \tag{46}
\]

将数值回代式（37）、式（38），两式均成立；这一步同时检验了滚动几何、符号和数值根的选择。

\subsection{六、关键洞察}\label{ux516dux5173ux952eux6d1eux5bdf-3}

\begin{enumerate}
\def\labelenumi{\arabic{enumi}.}
\tightlist
\item
  磁场不做功只保证机械能不变，并不保证机械角动量不变。
\item
  磁力矩是 \(r^2\) 的全微分，因而可以构造守恒的``机械角动量减位置项''。
\item
  只比较相对静止的两个时刻，可把多刚体瞬时运动统一为等效转动惯量，避开全过程微分方程。
\item
  第（3）问中小环自转与圆心公转大小相等而方向相反，它使小珠落在 \(AO\)
  上，从而使磁项只含 \(\cos\alpha\)。
\end{enumerate}

\subsection{七、易错点}\label{ux4e03ux6613ux9519ux70b9-3}

\begin{enumerate}
\def\labelenumi{\arabic{enumi}.}
\tightlist
\item
  直接写 \(L_z\) 守恒，漏掉外磁力矩。
\item
  在能量方程中加入``磁势能''；静磁场对点电荷不做功，本题机械动能本身守恒。
\item
  把均匀细环绕直径的转动惯量误用于垂直环面的转轴。
\item
  第（2）问求出代数根后不检查 \(r\leq2R\)。
\item
  内滚时把接触半径夹角 \(2\alpha\)
  当成小环相对大环的自转角，或把小环绕心惯量漏掉。
\item
  式（41）平方后不回代未平方的式（40），保留伪根。
\end{enumerate}

\subsection{八、课堂讲解建议}\label{ux516bux8bfeux5802ux8bb2ux89e3ux5efaux8bae-3}

先让学生计算 \(\bm F_B\cdot\bm v\)，确认能量守恒；再计算
\(\bm r\times\bm F_B\)，分别判断磁力对能量与角动量的影响。广义角动量式最好从坐标分量一步推到全微分，避免把它当作需要背诵的结论。

讲第三问时，先只画 \(A,O,O'\)
三点，独立完成余弦定理；再把小环自转惯量和小珠惯量逐项加入。每加一项都说明其物理对象、转轴和距离，学生就不容易把式（36）看成突然出现的代数式。

数值求根之后必须回到式（40）检查。课堂上可把被舍去的根也列出，让学生看到``数学方程的根''与``第一次再次相对静止的物理解''并不等价。

\subsection{九、变式训练}\label{ux4e5dux53d8ux5f0fux8badux7ec3-3}

\begin{enumerate}
\def\labelenumi{\arabic{enumi}.}
\tightlist
\item
  将磁场方向反转，重新写出广义角动量守恒式并讨论第（2）问是否仍有可达位置。
\item
  把小珠质量改为 \(M\)，推导式（21）---式（24）的对应形式。
\item
  把小环半径改为
  \(\lambda R\)，先求纯滚动的倍角关系，再建立末态等效转动惯量。
\item
  若铰点不在大环上而在距圆心 \(d\) 处，讨论哪些推导仍然成立。
\end{enumerate}

\needspace{3\baselineskip}
\textbf{变式解答：小环非平凡根的筛选}

令 \(c=\cos\alpha\)，取主解析的
\(\omega'_0=0.53qB/m\)。能量和正则角动量分别写成

\[
{5\over2}(0.53)^2=(c^2-3c+4.5)\left({m\omega_2\over qB}\right)^2,
\]

\[
{5\over2}(0.53)=(c^2-3c+4.5){m\omega_2\over qB}-{(1-c)^2\over2}.
\]

消去 \(\omega_2\)，得到

\[
\left[1+{(1-c)^2\over2.65}\right]^2
={2\over5}(c^2-3c+4.5).
\]

\(c=1\) 是初态根，不能当成完成滚动后的新位置。另一物理解为
\(c\simeq0.1335082\)，给出
\(\alpha\simeq82.3276^\circ\)、\(m\omega_2/(qB)\simeq0.4129901\)。

求根之后应同时回代两条守恒式，而不只回代消元所得的平方方程；平方过程可能保留角速度反号的额外分支。还要核对
\(-1\le c\le1\)，以及从初态连续运动能否首先到达此根。图中小珠始终在
\(AO\) 上，其距离为 \(R(1-c)\)，正则角动量的磁场项因此含
\((1-c)^2\)，并非只与小环自转角有关。

\needspace{3\baselineskip}
\textbf{小环消元方程的初态根与物理解}

把 \(u=1-c\) 代入原方程。由于 \(c^2-3c+4.5=2.5+u+u^2\)，得到

\[
\left(1+{u^2\over2.65}\right)^2
=1+0.4u+0.4u^2.
\]

移项并提出 \(u\)：

\[
u\left[{u^3\over2.65^2}
+\left({2\over2.65}-0.4\right)u-0.4\right]=0.
\]

\(u=0\) 对应初态。方括号在 \(u\ge0\) 的导数为

\[
{3u^2\over2.65^2}+{2\over2.65}-0.4>0,
\]

因此至多有一个非负根；它在 \(u=0\) 为负、\(u=2\) 为正，故恰有一个
\(0<u<2\) 的物理非平凡根。数值 \(u\simeq0.8664918\) 恢复主解析的
\(c\simeq0.1335082\)。

这项单调性证明排除了在允许余弦范围内漏掉另一条正 \(u\)
分支。仍须回到未平方的角动量式选择正向角速度，确认运动从初态连续到达该角度。小珠在初态重合点并非一般圆周轨道，不能把大环半径直接当作小珠到铰点的固定距离。

\newpage

\section{第 5 题}\label{ux7b2c-5-ux9898}

将一个线性分子简化为固定在同一直线上的三个点电荷，依次为
\(-q,2q,-q\)，相邻电荷间距均为 \(l\)。分子总质量为
\(m\)，绕质心且垂直于分子轴的转动惯量为 \(I=\lambda ml^2\)。以远处点电荷
\(Q\) 到分子质心的位移为 \(\bm r=r\hat{\bm r}\)，分子轴方向为
\(\hat{\bm n}\)，二者夹角为 \(\theta\)，且 \(r\gg l\)。

（1）保留最低非零阶，求该分子在点电荷 \(Q\) 的电场中受到的合力。

（2）求分子受到的合力矩，精确到相对量级 \((l/r)^2\)。

（3）若把分子置于圆形区域中心，区域内磁场均匀但随时间变化，且
\(\partial\bm B/\partial t=\bm b\)，求感生电场对分子产生的合力矩。用
\(\bm b,\hat{\bm n},l,q\) 表示。

\begin{figure}[H]
\centering
\includegraphics[width=0.86\linewidth,height=\textheight,keepaspectratio,alt={图5-1　线性三电荷分子的取向及其与远处点电荷的位置关系}]{../figures/question-05.pdf}
\caption*{图5-1　线性三电荷分子的取向及其与远处点电荷的位置关系}
\end{figure}

\subsection{一、题型定位}\label{ux4e00ux9898ux578bux5b9aux4f4d-4}

线性三电荷分子的总电荷和偶极矩均为零，远场作用从四极项开始。点电荷给出静电力与定向力矩，变化磁场则通过旋涡电场产生另一种力矩。

\subsection{二、知识点}\label{ux4e8cux77e5ux8bc6ux70b9-4}

\begin{enumerate}
\def\labelenumi{\arabic{enumi}.}
\tightlist
\item
  电势的泰勒展开与多极矩阶次。
\item
  零总电荷、零偶极矩和非零四极矩。
\item
  由势能求力：\(\bm F=-\bm\nabla U\)。
\item
  球坐标中径向与角向梯度。
\item
  由势能求取向力矩：\(\tau=-\partial U/\partial\theta\)。
\item
  法拉第电磁感应与对称规范下的感生电场。
\item
  向量三重积恒等式。
\end{enumerate}

\subsection{三、物理图像}\label{ux4e09ux7269ux7406ux56feux50cf-4}

若只从很远处观察三个电荷，最先看到的总电荷为零；再看正负电荷中心的分离，偶极矩也为零。只有电场曲率能够区分分子两端的
\(-q\) 与中心的 \(2q\)，所以作用强度比偶极效应多含一个 \(l/r\) 小量。

点电荷场既改变分子质心的运动，也试图改变分子轴取向。图 5-2 把合力分成
\(\hat{\bm r}\)
与角向分量：仅写径向分量会漏掉取向随空间方向变化所带来的横向力。感生电场部分则围绕
\(\bm b\) 形成闭合环流，只有 \(\bm b\)
在垂直分子轴方向的分量产生取向力矩。

\subsection{四、解题入口}\label{ux56dbux89e3ux9898ux5165ux53e3-4}

把分子质心作为展开点，三个电荷的位置写成

\[
\bm s_-=-l\hat{\bm n},\qquad
\bm s_0=\bm0,
\qquad
\bm s_+=l\hat{\bm n}.
\]

不要分别把三个库仑力展开到复杂的向量分量。先展开点电荷 \(Q\) 在
\(\bm r+\bm s\) 处的电势，求得分子势能，再对 \(\bm r\) 与 \(\theta\)
求导。这样总电荷项和偶极项会自动抵消，也不易漏项。

\subsection{五、解析}\label{ux4e94ux89e3ux6790-4}

\begin{figure}[H]
\centering
\includegraphics[width=0.9\linewidth,height=\textheight,keepaspectratio,alt={图5-2　四极势能产生的合力、取向力矩及感生电场力矩}]{../figures/solution-05.pdf}
\caption*{图5-2　四极势能产生的合力、取向力矩及感生电场力矩}
\end{figure}

图5-2中，左图同时固定 \(\hat{\bm r}\)、分子轴与 \(\hat{\bm\theta}\)
的方向，力的分量按此取符号。右图给出垂直磁场变化率下两端相反的力，说明合力为零仍可有力矩。

\needspace{3\baselineskip}
\textbf{（一）确认最低非零多极阶次}

三个电荷的总电荷为

\[
Q_{\rm mol}=(-q)+2q+(-q)=0. \tag{1}
\]

相对于质心的电偶极矩为

\[
\begin{aligned}
\bm p
&=(-q)(-l\hat{\bm n})+(2q)\bm0+(-q)(l\hat{\bm n})\\
&=\bm0. 
\end{aligned}
\tag{2}
\]

故单极项和偶极项均消失，最低非零项是二阶空间展开，也就是电四极项。

\needspace{3\baselineskip}
\textbf{（二）势能的二阶展开}

点电荷 \(Q\) 在位置 \(\bm r+\bm s\) 处产生的电势能因子为

\[
\frac1{|\bm r+\bm s|}
=\frac1r\left(1+\frac{2\bm r\cdot\bm s+s^2}{r^2}\right)^{-1/2}. \tag{3}
\]

使用二项式展开

\[
(1+x)^{-1/2}=1-\frac12x+\frac38x^2+O(x^3), \tag{4}
\]

并只保留到 \(s^2/r^2\)，得到

\[
\frac1{|\bm r+\bm s|}
=\frac1r-\frac{\hat{\bm r}\cdot\bm s}{r^2}
+\frac{3(\hat{\bm r}\cdot\bm s)^2-s^2}{2r^3}
+O\!\left(\frac{s^3}{r^4}\right). \tag{5}
\]

对中心电荷 \(2q\)，\(\bm s=0\)；对两端电荷，\(s=l\) 且

\[
\hat{\bm r}\cdot\bm s_\pm=\pm l\cos\theta. \tag{6}
\]

分子在 \(Q\) 的电场中的总势能为

\[
U=k_eQ\sum_i\frac{q_i}{|\bm r+\bm s_i|},
\qquad k_e=\frac1{4\pi\varepsilon_0}. \tag{7}
\]

式（5）的零阶项乘以 \(\sum q_i=0\)，一阶项乘以
\(\sum q_i\bm s_i=0\)。二阶项只由两个端电荷贡献：

\[
\begin{aligned}
U
&=k_eQ\cdot2(-q)
\frac{3l^2\cos^2\theta-l^2}{2r^3}\\
&=\frac{k_eQql^2}{r^3}(1-3\cos^2\theta). 
\end{aligned}
\tag{8}
\]

令

\[
K=k_eQql^2=\frac{Qql^2}{4\pi\varepsilon_0}, \tag{9}
\]

则

\[
\boxed{U=\frac K{r^3}(1-3\cos^2\theta)}. \tag{10}
\]

量纲检查：\(k_eQq/r\) 是能量，乘 \(l^2/r^2\) 后仍为能量；并且
\(U\propto r^{-3}\)，符合四极势能的远场阶次。

\needspace{3\baselineskip}
\textbf{（三）由势能求完整合力}

令 \(c=\cos\theta\)。在 \(\hat{\bm r}\) 与分子轴张成的平面内，梯度为

\[
\bm\nabla U
=\frac{\partial U}{\partial r}\hat{\bm r}
+\frac1r\frac{\partial U}{\partial\theta}\hat{\bm\theta}. \tag{11}
\]

先算两个偏导：

\[
\frac{\partial U}{\partial r}
=-\frac{3K}{r^4}(1-3c^2), \tag{12}
\]

\[
\frac{\partial U}{\partial\theta}
=\frac{6K}{r^3}\sin\theta\cos\theta. \tag{13}
\]

所以

\[
\begin{aligned}
\bm F=-\bm\nabla U
&=\frac{3K}{r^4}(1-3c^2)\hat{\bm r}
-\frac{6K}{r^4}\sin\theta\cos\theta\hat{\bm\theta}. 
\end{aligned}
\tag{14}
\]

即

\[
\boxed{
\bm F
=\frac{3Qql^2}{4\pi\varepsilon_0r^4}
\left[(1-3\cos^2\theta)\hat{\bm r}
-2\sin\theta\cos\theta\hat{\bm\theta}\right]}. \tag{15}
\]

式（15）是完整合力。若题目只问沿 \(\bm r\) 的分量，则

\[
\boxed{F_r=\frac{3Qql^2}{4\pi\varepsilon_0r^4}(1-3\cos^2\theta)}. \tag{16}
\]

还可以把式（14）写成不依赖角向基矢的形式：

\[
\boxed{
\bm F=\frac{3K}{r^4}
\left[(1-5c^2)\hat{\bm r}+2c\hat{\bm n}\right]}. \tag{17}
\]

式（17）展开到 \(\hat{\bm r},\hat{\bm\theta}\) 后与式（14）一致。特殊角
\(c^2=1/3\)
只会使径向分量为零，横向分量一般仍不为零，因此不能说总力必为零。

\needspace{3\baselineskip}
\textbf{（四）分子受到的取向力矩}

分子轴发生微小转角 \(\mathrm d\theta\) 时，取向力矩在垂直于
\(\hat{\bm r},\hat{\bm n}\) 平面的分量满足

\[
\tau_\theta=-\frac{\partial U}{\partial\theta}. \tag{18}
\]

由式（13）：

\[
\boxed{
\tau_\theta=-\frac{6K}{r^3}\sin\theta\cos\theta}. \tag{19}
\]

用向量表示，注意 \(|\hat{\bm r}\times\hat{\bm n}|=\sin\theta\)：

\[
\boxed{
\bm\tau
=-\frac{6K}{r^3}\cos\theta
(\hat{\bm r}\times\hat{\bm n})}. \tag{20}
\]

当 \(Qq>0\) 时，\(\theta=0\)
附近势能最低，力矩使分子轴趋向径向；\(\theta=\pi/2\)
是另一平衡取向，但稳定性相反。将 \(U\) 对 \(\theta\)
做二阶导数可验证这一判断。

\needspace{3\baselineskip}
\textbf{（五）时变磁场产生的感生电场}

区域内 \(\partial\bm B/\partial t=\bm b\) 为空间常量。由法拉第定律

\[
\bm\nabla\times\bm E=-\bm b. \tag{21}
\]

在以分子质心为原点的圆形对称区域中，可取

\[
\boxed{\bm E(\bm x)=-\frac12\bm b\times\bm x}. \tag{22}
\]

因为对常向量 \(\bm b\) 有

\[
\bm\nabla\times(\bm b\times\bm x)=2\bm b, \tag{23}
\]

故式（22）确实满足式（21）。中心电荷 \(2q\) 位于
\(\bm x=0\)，所受感生电场力为零。两个端电荷的位置分别为

\[
\bm x_\pm=\pm l\hat{\bm n}. \tag{24}
\]

端电荷均为 \(-q\)，它们所受力为

\[
\bm F_\pm=(-q)\bm E(\bm x_\pm)
=\frac q2\bm b\times\bm x_\pm. \tag{25}
\]

每个端电荷对质心的力矩为 \(\bm x_\pm\times\bm F_\pm\)。两端相加：

\[
\begin{aligned}
\bm\tau
&=\sum_{\pm}\bm x_\pm\times
\left(\frac q2\bm b\times\bm x_\pm\right)\\
&=q l^2\hat{\bm n}\times(\bm b\times\hat{\bm n}). 
\end{aligned}
\tag{26}
\]

利用向量三重积

\[
\hat{\bm n}\times(\bm b\times\hat{\bm n})
=\bm b(\hat{\bm n}\cdot\hat{\bm n})
-\hat{\bm n}(\hat{\bm n}\cdot\bm b), \tag{27}
\]

得到

\[
\boxed{
\bm\tau
=ql^2\left[\bm b-\hat{\bm n}
(\hat{\bm n}\cdot\bm b)\right]}. \tag{28}
\]

等价地，

\[
\boxed{
\bm\tau=-ql^2\hat{\bm n}\times
(\hat{\bm n}\times\bm b)}. \tag{29}
\]

若
\(\hat{\bm n}\parallel\bm b\)，两个端点都位于磁场变化的轴线上，感生电场各自为零，式（28）给出零；若
\(\hat{\bm n}\perp\bm b\)，力矩大小达到 \(ql^2|\bm b|\)。

\subsection{六、关键洞察}\label{ux516dux5173ux952eux6d1eux5bdf-4}

\begin{enumerate}
\def\labelenumi{\arabic{enumi}.}
\tightlist
\item
  总电荷与偶极矩同时为零，最低非零相互作用必从二阶展开开始。
\item
  先求势能再求力，比直接展开三组库仑力更短且更不易漏掉横向分量。
\item
  ``径向力为零''不等于``合力为零''，因为分子取向使势能也依赖角坐标。
\item
  感生电场对位置是线性的，正好把端电荷间距的二次量 \(l^2\) 带入力矩。
\end{enumerate}

\subsection{七、易错点}\label{ux4e03ux6613ux9519ux70b9-4}

\begin{enumerate}
\def\labelenumi{\arabic{enumi}.}
\tightlist
\item
  把该体系误当成电偶极子，写出 \(r^{-3}\) 的偶极力。
\item
  泰勒展开只保留一阶，得到零后误判分子完全不受力。
\item
  由 \(U(r,\theta)\) 求力时只算
  \(-\partial U/\partial r\)，漏掉角向分量。
\item
  力矩中漏写一个 \(\sin\theta\) 或把 \(r^{-3}\) 写成 \(r^{-4}\)。
\item
  感生电场写成 \(-\bm b\times\bm x\)，漏掉系数 \(1/2\)。
\item
  向量三重积次序颠倒，导致式（28）整体符号错误。
\end{enumerate}

\subsection{八、课堂讲解建议}\label{ux516bux8bfeux5802ux8bb2ux89e3ux5efaux8bae-4}

先让学生只计算 \(\sum q_i\) 和
\(\sum q_i\bm s_i\)，据此预测答案的距离幂次，再进入泰勒展开。这样每一项为什么消失都有物理原因，而不是代数巧合。

讲合力时可选 \(\theta=0,\pi/2\) 和 \(\cos^2\theta=1/3\)
三个特例检查式（15）。尤其第三个特例能直接揭示``只给径向分量''的不足。

第三问先用圆周积分验证
\(E_\varphi=-\tfrac12br\)，再写成向量式（22）。最后用平行与垂直两个极限检查力矩方向和大小。

\subsection{九、变式训练}\label{ux4e5dux53d8ux5f0fux8badux7ec3-4}

\begin{enumerate}
\def\labelenumi{\arabic{enumi}.}
\tightlist
\item
  将三个电荷改为 \(q,-2q,q\)，讨论势能、力和力矩的符号变化。
\item
  给分子增加一个很小的永久偶极矩
  \(\bm p\)，比较偶极项与四极项在多远处交叉。
\item
  求 \(Qq>0\) 时两个平衡取向的稳定性及小角振动频率。
\item
  若 \(\bm b\) 与分子轴夹角为 \(\chi\)，求感生力矩大小并讨论其趋向。
\end{enumerate}

\needspace{3\baselineskip}
\textbf{变式解答：分子轴平行与垂直径向}

保持 \(K=k_eQql^2\)，静电能为
\(U=K(1-3\cos^2\theta)/r^3\)。当分子轴平行径向，\(\theta=0\)，有

\[
U=-{2K\over r^3},\qquad F_r=-{6K\over r^4},\qquad F_\theta=0.
\]

当分子轴垂直径向，\(\theta=\pi/2\)，有

\[
U={K\over r^3},\qquad F_r={3K\over r^4},\qquad F_\theta=0.
\]

两种方向都没有静电力矩，但能量曲率相反：

\[
{\partial^2U\over\partial\theta^2}={6K\over r^3}\cos2\theta.
\]

\(K>0\) 时平行方向稳定、垂直方向不稳定；\(K<0\)
时互换。零力矩只说明平衡，并不保证稳定。

若磁场变化率 \(\bm b\)
平行分子轴，各电荷位置均在该轴线上，\(\bm b\times\bm r_i=0\)，每个电荷上的旋涡电场力都为零。若
\(\bm b\perp\hat{\bm n}\)，两端力反向、合力为零、力矩为
\(ql^2\bm b\)。这与图中方向一致，并区分了``合力为零''和``各电荷均不受力''两种情形。

\newpage

\section{第 6 题}\label{ux7b2c-6-ux9898}

浅水波的波速近似为 \(v=\sqrt{gh}\)，其中 \(h\) 是水深。取岸线为 \(x\)
轴，海面位于 \(y>0\)。水深随 \(y\) 变化：

\[
h(y)=\frac{4h_0}{5-y^2/y_0^2},
\]

其中 \(h_0,y_0\) 为常量。一艘船沿 \(y=y_0\) 以速度 \(v_b=2\sqrt{gh_0}\)
向 \(+x\) 方向匀速航行，\(t=0\) 时船位于
\(A(0,y_0)\)。岸边观察者位于原点 \(O\)。

（1）船在 \(A\) 处激发的一束水波射线到达岸边 \(B\)，记 \(x_1=|OB|\)，求
\(x_1/y_0\)。

（2）观察者第一次接收到船激发的水波的时刻为 \(t_r\)，求
\(t_r\sqrt{gh_0}/y_0\)。

（3）\(t=0\) 时水面上的波前可写为 \(y=f(x)\)。若 \(f(x_s)=y_0/2\)，求
\(x_s/y_0\)。

\begin{figure}[H]
\centering
\includegraphics[width=0.88\linewidth,height=\textheight,keepaspectratio,alt={图6-1　变水深海域中的运动波源、岸线与弯曲水波射线}]{../figures/question-06.pdf}
\caption*{图6-1　变水深海域中的运动波源、岸线与弯曲水波射线}
\end{figure}

\subsection{一、题型定位}\label{ux4e00ux9898ux578bux5b9aux4f4d-5}

浅水波速随离岸距离变化，射线会折弯。船的运动决定初始波矢方向；射线表示一束波的传播路径，波前则由不同发射时刻的同相位点组成。

\subsection{二、知识点}\label{ux4e8cux77e5ux8bc6ux70b9-5}

\begin{enumerate}
\def\labelenumi{\arabic{enumi}.}
\tightlist
\item
  浅水波速 \(v=\sqrt{gh}\)。
\item
  分层介质中的 Snell 不变量 \(\cos\psi/v\)。
\item
  射线轨迹微分方程与变量代换。
\item
  沿弯曲射线的真实传播时间积分。
\item
  匀速运动波源的发射位置与发射时刻。
\item
  波前与射线的正交关系。
\item
  参数方程消元及数值核验。
\end{enumerate}

\subsection{三、物理图像}\label{ux4e09ux7269ux7406ux56feux50cf-5}

本题水深从 \(y=y_0\) 向岸边逐渐减小，波速随之改变。由于介质只在 \(y\)
方向不均匀，波矢的 \(x\) 分量保持不变，射线会连续弯折。选择能从 \(A\)
到达岸线的射线后，其轨迹恰好是一段余弦曲线。

船在 \(t<0\) 时位于 \(x<0\)。观察者在原点第一次听到的并不是 \(t=0\) 才从
\(A\) 发出的波，而是船更早在 \(A'\) 发出的波。\(t=0\)
的波前也由这一族历史波束共同组成：不同点对应不同的发射位置、传播时长和射线参数。

\subsection{四、解题入口}\label{ux56dbux89e3ux9898ux5165ux53e3-5}

先记

\[
v_0=\sqrt{gh_0},\qquad Y=\frac y{y_0},\qquad X=\frac x{y_0}. 
\]

把波速化成无量纲形式。设射线速度与 \(+x\) 方向夹角为
\(\psi\)，由运动波源的包络条件 \(v_b(k_x/\omega)=1\)
选出能形成船波波前的射线。取指向岸边的分支，在 \(A\) 点有，使
\(v_x=v_0/2\)、\(v_y=-\sqrt3v_0/2\)。利用 \(\cos\psi/v\) 不变建立
\(\mathrm dY/\mathrm dX\) 的方程。

第一问得到射线的水平位移后，第二问可直接平移这条射线，使其从船早先的位置到达
\(O\)。第三问则引入参数 \(\phi\)
表示每个波前点所在射线的局部相位角，分别写出 \(X(\phi),Y(\phi)\)。

\subsection{五、解析}\label{ux4e94ux89e3ux6790-5}

\begin{figure}[H]
\centering
\includegraphics[width=0.9\linewidth,height=\textheight,keepaspectratio,alt={图6-2　单条射线路径与同一时刻波前的区别}]{../figures/solution-06.pdf}
\caption*{图6-2　单条射线路径与同一时刻波前的区别}
\end{figure}

图6-2中，左图只画有效域内的 \(A\to B\) 射线，右图画同一时刻的波前。将
\(Y=1/2\) 代入参数式得到岸前波面的位置；两条曲线相遇处的斜率乘积为
\(-1\)。

\needspace{3\baselineskip}
\textbf{（一）波速分布与射线不变量}

由题给水深

\[
h(y)=\frac{4h_0}{5-Y^2}, \tag{1}
\]

得到

\[
\boxed{v(Y)=\frac{2v_0}{\sqrt{5-Y^2}}}. \tag{2}
\]

在 \(Y=1\) 的航线上，\(v=v_0\)。介质不显含 \(x\)，故波矢的 \(x\) 分量
\(k_x\) 守恒。对无色散波
\(\omega=vk\)，频率沿静止介质中的同一条射线也不变，所以

\[
\frac{k_x}{\omega}=\frac{\cos\psi}{v}=\text{常量}. \tag{3}
\]

需要先说明为何初始方向满足
\(\cos\psi_A=1/2\)。对持续匀速运动的波源，稳态波型满足
\(\omega=\bm k\cdot\bm v_b\)，而浅水无色散关系为 \(\omega=v k\)。所以

\[
{\cos\psi\over v}={1\over v_b}={1\over2v_0}.
\]

等价地，首次到达的包络满足对发射时刻的驻值条件 \(1-v_b p_x=0\)，其中
\(p_x=k_x/\omega\) 是守恒的沿岸慢度。本题不是任选一条从 \(A\)
发出的射线；若不给出这一包络条件，第一问会有一族可能的岸边落点。在 \(A\)
点，该射线的速度分量为

\[
v_x(A)=\frac{v_0}{2},\qquad
v_y(A)=-\frac{\sqrt3v_0}{2}. \tag{4}
\]

因此 \(\cos\psi_A=1/2\)，式（3）给出

\[
\frac{\cos\psi}{v}=\frac1{2v_0}. \tag{5}
\]

代入式（2）：

\[
\cos\psi=\frac{v}{2v_0}
=\frac1{\sqrt{5-Y^2}}. \tag{6}
\]

由于射线向岸边传播，取 \(\sin\psi<0\)：

\[
\sin\psi
=-\sqrt{1-\frac1{5-Y^2}}
=-\sqrt{\frac{4-Y^2}{5-Y^2}}. \tag{7}
\]

于是射线斜率

\[
\frac{\mathrm dy}{\mathrm dx}
=\tan\psi
=-\sqrt{4-Y^2}. \tag{8}
\]

注意式（8）右端无量纲；因为 \(X=x/y_0,Y=y/y_0\)，也有

\[
\boxed{\frac{\mathrm dY}{\mathrm dX}=-\sqrt{4-Y^2}}. \tag{9}
\]

\needspace{3\baselineskip}
\textbf{（二）从 \(A\) 到岸线的射线路径}

对式（9）分离变量：

\[
\frac{\mathrm dY}{\sqrt{4-Y^2}}=-\mathrm dX. \tag{10}
\]

积分得到

\[
\arcsin\frac Y2=-X+C. \tag{11}
\]

在 \(A\) 点 \(X=0,Y=1\)，故

\[
C=\arcsin\frac12=\frac\pi6. \tag{12}
\]

于是

\[
\arcsin\frac Y2=\frac\pi6-X. \tag{13}
\]

把正弦形式改写成便于描述本段轨迹的余弦形式：

\[
\boxed{Y=2\cos\left(X+\frac\pi3\right)}. \tag{14}
\]

即

\[
\boxed{y=2y_0\cos\left(\frac x{y_0}+\frac\pi3\right)}. \tag{15}
\]

到达岸线时 \(Y=0\)。本段轨迹中括号内角度从 \(\pi/3\) 增到
\(\pi/2\)，因此

\[
X_1+\frac\pi3=\frac\pi2,\qquad X_1=\frac\pi6. \tag{16}
\]

所以

\[
\boxed{\frac{x_1}{y_0}=\frac\pi6}. \tag{17}
\]

\needspace{3\baselineskip}
\textbf{（三）单条射线的真实传播时间}

不能把式（14）中的参数 \(X\) 当成匀速时间。射线的水平速度为

\[
\frac{\mathrm dx}{\mathrm dt}=v\cos\psi. \tag{18}
\]

由式（2）、式（6）：

\[
v\cos\psi
=\frac{2v_0}{\sqrt{5-Y^2}}
\frac1{\sqrt{5-Y^2}}
=\frac{2v_0}{5-Y^2}. \tag{19}
\]

故

\[
\mathrm dt=\frac{y_0}{2v_0}(5-Y^2)\,\mathrm dX. \tag{20}
\]

令

\[
\phi=X+\frac\pi3,\qquad Y=2\cos\phi. \tag{21}
\]

从 \(A\) 到岸边时 \(\phi:\pi/3\to\pi/2\)，且
\(\mathrm dX=\mathrm d\phi\)。所以传播时间

\[
\begin{aligned}
t_{\rm prop}
&=\frac{y_0}{2v_0}
\int_{\pi/3}^{\pi/2}(5-4\cos^2\phi)\,\mathrm d\phi\\
&=\frac{y_0}{2v_0}
\int_{\pi/3}^{\pi/2}(3-2\cos2\phi)\,\mathrm d\phi\\
&=\frac{y_0}{2v_0}
\left[3\phi-\sin2\phi\right]_{\pi/3}^{\pi/2}\\
&=\frac{y_0}{2v_0}
\left(\frac\pi2+\frac{\sqrt3}{2}\right)\\
&=\boxed{\frac{y_0}{4v_0}(\pi+\sqrt3)}. 
\end{aligned}
\tag{22}
\]

上式的第二行使用
\(\cos^2\phi=(1+\cos2\phi)/2\)。量纲为长度除以速度，正确。

\needspace{3\baselineskip}
\textbf{（四）观察者第一次接收到水波的时刻}

从航线上的任一点发出、采用同样初始方向的射线，只是式（15）在 \(x\)
方向平移，水平传播距离始终为 \(x_1=\pi y_0/6\)。要使射线最终到达原点
\(O\)，发射点 \(A'\) 的横坐标必须为

\[
x_{A'}=-x_1=-\frac\pi6y_0. \tag{23}
\]

船速为 \(v_b=2v_0\)，且 \(t=0\) 时经过 \(x=0\)，所以对应发射时刻

\[
t_e=\frac{x_{A'}}{v_b}
=-\frac{\pi y_0}{12v_0}. \tag{24}
\]

接收时刻等于发射时刻加传播时长：

\[
\begin{aligned}
t_r
&=t_e+t_{\rm prop}\\
&=-\frac{\pi y_0}{12v_0}
+\frac{y_0}{4v_0}(\pi+\sqrt3)\\
&=\frac{y_0}{v_0}
\left(\frac\pi6+\frac{\sqrt3}{4}\right). 
\end{aligned}
\tag{25}
\]

因此

\[
\boxed{
\frac{t_r\sqrt{gh_0}}{y_0}
=\frac\pi6+\frac{\sqrt3}{4}}. \tag{26}
\]

这里 \(t_r>0\)，说明虽然波在 \(t<0\)
已发出，但传播耗时更长，观察者在船经过 \(A\) 之后才第一次收到。

\needspace{3\baselineskip}
\textbf{（五）\(t=0\) 时的波前参数方程}

波前上的任一点 \(P\) 来自某个早于 \(0\) 的发射时刻。用 \(\phi\)
表示该点所在射线从发射点走到 \(P\) 后，式（21）中的终值。发射点处
\(\phi_0=\pi/3\)，故

\[
\boxed{Y=2\cos\phi},\qquad
\frac\pi3\leq\phi\leq\frac\pi2. \tag{27}
\]

从发射点到 \(P\) 的无量纲水平位移为

\[
\Delta X=\phi-\frac\pi3. \tag{28}
\]

相应传播时间由式（20）积分到一般上限 \(\phi\)：

\[
\begin{aligned}
\frac{v_0\Delta t}{y_0}
&=\frac12\int_{\pi/3}^{\phi}(5-4\cos^2u)\,\mathrm du\\
&=\frac12\left[3\left(\phi-\frac\pi3\right)
-\sin2\phi+\frac{\sqrt3}{2}\right]. 
\end{aligned}
\tag{29}
\]

若该波在 \(t_e=-\Delta t\) 发射，船的无量纲发射位置为

\[
X_e=\frac{v_bt_e}{y_0}
=-2\frac{v_0\Delta t}{y_0}. \tag{30}
\]

波前点的坐标等于发射位置加射线位移：

\[
X=X_e+\Delta X. \tag{31}
\]

把式（28）---式（30）代入并逐项合并：

\[
\begin{aligned}
X
&=-\left[3\left(\phi-\frac\pi3\right)
-\sin2\phi+\frac{\sqrt3}{2}\right]
+\left(\phi-\frac\pi3\right)\\
&=-2\left(\phi-\frac\pi3\right)
+\sin2\phi-\frac{\sqrt3}{2}. 
\end{aligned}
\tag{32}
\]

因此 \(t=0\) 波前的参数方程为

\[
\boxed{
\begin{aligned}
\frac y{y_0}&=2\cos\phi,\\
\frac x{y_0}&=-2\left(\phi-\frac\pi3\right)
+\sin2\phi-\frac{\sqrt3}{2},
\end{aligned}
\quad
\frac\pi3\leq\phi\leq\frac\pi2.} \tag{33}
\]

\needspace{3\baselineskip}
\textbf{（六）求 \(y=y_0/2\) 处的横坐标}

由第一式

\[
2\cos\phi=\frac12,\qquad
\cos\phi=\frac14. \tag{34}
\]

在指定区间内

\[
\phi=\arccos\frac14. \tag{35}
\]

又因为

\[
\sin\phi=\frac{\sqrt{15}}4,\qquad
\sin2\phi=2\sin\phi\cos\phi=\frac{\sqrt{15}}8, \tag{36}
\]

代入式（32）：

\[
\boxed{
\frac{x_s}{y_0}
=-2\left(\arccos\frac14-\frac\pi3\right)
+\frac{\sqrt{15}}8-\frac{\sqrt3}{2}}. \tag{37}
\]

数值为

\[
\boxed{\frac{x_s}{y_0}=-0.9237395264\approx-0.9237}. \tag{38}
\]

负号与图 6-2 一致：\(t=0\) 时到达 \(y=y_0/2\) 的这部分波前来自船早先位于
\(x<0\) 的时刻。

\needspace{3\baselineskip}
\textbf{波前与射线的正交校验}

由式（33），有
\(\mathrm dX/\mathrm d\phi=-4\sin^2\phi\)、\(\mathrm dY/\mathrm d\phi=-2\sin\phi\)，所以波前切线斜率为
\(1/(2\sin\phi)\)。同一位置的射线斜率由式（9）为
\(-2\sin\phi\)，两者乘积为
\(-1\)。这验证了式（33）确实描述包络波前，而不是另一条射线；图6-2同时标出这两条方向。

\subsection{六、关键洞察}\label{ux516dux5173ux952eux6d1eux5bdf-5}

\begin{enumerate}
\def\labelenumi{\arabic{enumi}.}
\tightlist
\item
  介质沿 \(x\) 方向平移对称，守恒的是横向慢度，而不是水平速度。
\item
  射线方程描述能量传播路径；波前描述同一时刻的等相位点，两者不是同一条曲线。
\item
  第一次接收必须允许发射时刻为负，不能只研究 \(t\geq0\) 后的波。
\item
  式（14）的参数 \(X\) 是几何坐标，传播速度不均匀，求时间必须重新积分。
\end{enumerate}

\subsection{七、易错点}\label{ux4e03ux6613ux9519ux70b9-5}

\begin{enumerate}
\def\labelenumi{\arabic{enumi}.}
\tightlist
\item
  把 \(v\sin\psi\) 或 \(v\cos\psi\) 当作分层介质中的守恒量。
\item
  在平方根中选错 \(\sin\psi\) 的符号，使射线朝远离岸线方向传播。
\item
  用直线距离除以 \(v_0\) 估算传播时间，忽略路径弯曲和变速。
\item
  把 \(A\) 点发出的射线到达 \(B\) 的时刻当成观察者在 \(O\)
  的首次接收时刻。
\item
  第三问直接使用单条射线式（15），没有加入波源在发射时刻的位置。
\item
  计算 \(\sin2\phi\) 时漏掉系数 \(2\)。
\end{enumerate}

\subsection{八、课堂讲解建议}\label{ux516bux8bfeux5802ux8bb2ux89e3ux5efaux8bae-5}

先在均匀介质中区分``射线''和``圆形波前''，再转入分层介质。学生一旦接受二者正交但不相同，第三问的参数化就不显得突兀。

第一问可让学生用两种积分变量完成：直接对 \(Y\) 积分，或设
\(Y=2\cos\phi\)。第二种变量会自然延续到传播时间和波前，体现选参的统一性。

每得到一个结果都做图像检查：射线在 \(A\) 的初始斜率应为
\(-\sqrt3\)；到岸时 \(Y=0\)；波前点 \(x_s\)
应位于原点左侧。图像与符号一致是本题最有效的错误筛查。

\subsection{九、变式训练}\label{ux4e5dux53d8ux5f0fux8badux7ec3-5}

\begin{enumerate}
\def\labelenumi{\arabic{enumi}.}
\tightlist
\item
  若船速改为 \(v_b=\eta v_0\)，推导首次接收时刻关于 \(\eta\) 的表达式。
\item
  若射线在 \(A\) 点与 \(+x\) 方向夹角为一般角
  \(\psi_0\)，建立轨迹方程并讨论能否到岸。
\item
  求 \(t=t_*\) 时刻的波前参数方程，说明它相对式（33）如何变化。
\item
  用式（33）计算波前切线，并验证其与相应射线方向正交。
\end{enumerate}

\needspace{3\baselineskip}
\textbf{变式解答：用射线斜率核对波前}

射线为 \(Y=2\cos(X+\pi/3)\)，所以

\[
{\mathrm dY\over\mathrm dX}=-2\sin(X+\pi/3).
\]

同一时刻的波前参数为

\[
X=-2(\phi-\pi/3)+\sin2\phi-\sqrt3/2,\qquad Y=2\cos\phi.
\]

对参数求导：

\[
{\mathrm dX\over\mathrm d\phi}=-4\sin^2\phi,\qquad
{\mathrm dY\over\mathrm d\phi}=-2\sin\phi,
\]

故波前斜率为 \(1/(2\sin\phi)\)，与当地射线斜率的乘积为
\(-1\)。两图应使用相同的 \(X,Y\)
物理单位；若绘图横纵比例不同，视觉上的角度会变，仍要由参数式判定正交。

波前在 \(Y=1/2\) 处有 \(\phi=\arccos(1/4)\)，代入得到
\(X\simeq-0.92374\)。这不是主射线在 \(Y=1/2\)
处的横坐标，因为波前包含更早发出的同相位波。

到岸传播时间也须沿弯曲射线积累：\(\mathrm ds/\mathrm dX=\sqrt{5-Y^2}\)、\(v/v_0=2/\sqrt{5-Y^2}\)，所以时间积分被积函数是
\((5-Y^2)/2\)。少掉这个因子 2，会把原题接收时间减半。

\newpage

\section{第 7 题}\label{ux7b2c-7-ux9898}

一根截面均匀、足够细的长玻璃管弯成图示形状。左端封闭，右端开口通大气。各段长度为

\[
l_1=100\,\mathrm{cm},\qquad
l_2=5\,\mathrm{cm},\qquad
l_3=50\,\mathrm{cm}.
\]

在左侧竖直管中、\(A\) 点下方有长度

\[
l_4=10\,\mathrm{cm}
\]

的水银柱。忽略表面张力、弯角处体积以及水银与管壁的阻力。大气压为
\(p_0=76\,\mathrm{cmHg}\)，初始环境温度为
\(T_0=300\,\mathrm K\)。管壁导热足够快，封闭气体各处始终均匀且与环境同温。

（1）判断初始平衡是否稳定。

（2）缓慢升温到多少温度 \(T_1\) 时，水银开始自发快速流动？

（3）承接第（2）问，当水银柱后端第一次到达 \(A\) 点时，求水银柱速度
\(v_1\)。

（4）此时求 \(B\) 点处水银内部的压强。

（5）水银继续运动，是否会从右端管口溢出？一旦溢出便不再讨论后续回流。

（6）待水银缓慢停下并再次达到平衡后，继续缓慢升温。若要使水银全部流出，至少应升温到多少？

\begin{figure}[H]
\centering
\includegraphics[width=0.64\linewidth,height=\textheight,keepaspectratio,alt={图7-1　封闭气体、水银柱与弯折细管各段长度的定义}]{../figures/question-07.pdf}
\caption*{图7-1　封闭气体、水银柱与弯折细管各段长度的定义}
\end{figure}

\subsection{一、题型定位}\label{ux4e00ux9898ux578bux5b9aux4f4d-6}

水银柱越过弯角后，两端高度差随位置分段改变。气体压力与所需静止压力的斜率决定初态稳定性，运动速度需把气体做功、大气做功和水银势能逐项相加。

\subsection{二、知识点}\label{ux4e8cux77e5ux8bc6ux70b9-6}

\begin{enumerate}
\def\labelenumi{\arabic{enumi}.}
\tightlist
\item
  等温理想气体状态方程与压强---体积曲线。
\item
  平衡点的微扰稳定性判据。
\item
  等温膨胀功 \(W=p_iV_i\ln(V_f/V_i)\)。
\item
  大气压力做功与液柱重力势能。
\item
  弯管中液柱沿管方向的一维动力学。
\item
  用机械能判定某构形是否可达。
\item
  临界全部流出时``最后一滴''的准静态条件。
\end{enumerate}

\subsection{三、物理图像}\label{ux4e09ux7269ux7406ux56feux50cf-6}

初始水银柱竖直位于 \(A\) 下方，左侧封闭气体压强比大气压高
\(10\,\mathrm{cmHg}\)。若水银向右前进一小段，封闭气体膨胀、压强降低；与此同时水银两端高度差也减小，维持平衡所需的压强降低得更快。因此气体仍显得``偏强''，继续推动水银，初态是不稳定平衡。

当水银后端到达 \(A\)
时，它的后半段位于水平段，前半段位于右侧向下的竖直段，质心比初态上升。再往后运动，水银整体逐渐下降，重力会释放能量，但封闭气体膨胀和大气反压共同决定它能否抵达开口。

\subsection{四、解题入口}\label{ux56dbux89e3ux9898ux5165ux53e3-6}

统一用``厘米水银柱''作为压强头。令水银从初始位置沿管向右移动距离
\(x\)，则封闭气柱长度从 \(90\,\mathrm{cm}\) 增为
\(90+x\)。先分别写出实际气体压强 \(p_g(x)\) 与若想静止所需的平衡压强
\(p_{\rm eq}(x)\)，比较它们在 \(x=0\) 附近的一阶项。

快速运动阶段选择整根水银为研究对象。气体功、大气功和重力势能必须使用同一初末构形。第（4）问再把水银分成
\(A\) 到 \(B\) 的水平半柱和 \(B\)
以下的竖直半柱，两段具有相同的沿管加速度大小。

\subsection{五、解析}\label{ux4e94ux89e3ux6790-6}

\begin{figure}[H]
\centering
\includegraphics[width=0.9\linewidth,height=\textheight,keepaspectratio,alt={图7-2　沿管坐标、气体压强曲线与平衡压强曲线}]{../figures/solution-07.pdf}
\caption*{图7-2　沿管坐标、气体压强曲线与平衡压强曲线}
\end{figure}

图7-2中，左图明确封闭端与水银初态，右图按水银两端所在管段绘制分段静止压力。下面的功积分使用同一个气柱长度
\(q\)；理想无阻力运动与长期弱阻尼静止须分开。

\needspace{3\baselineskip}
\textbf{（一）初态压强和气柱长度}

初态封闭气柱从左端延伸到水银柱下端，长度为

\[
q_0=l_1-l_4=100-10=90\,\mathrm{cm}. \tag{1}
\]

水银柱上端与经右侧开口相通的大气处于同一压强链路。初态竖直水银柱长
\(l_4\)，故封闭气体压强为

\[
p_i=p_0+\rho g l_4. \tag{2}
\]

以厘米水银柱为压强单位时，\(\rho gl_4\) 的数值就是
\(10\,\mathrm{cmHg}\)，所以

\[
\boxed{p_i=86\,\mathrm{cmHg}}. \tag{3}
\]

\needspace{3\baselineskip}
\textbf{（二）第（1）问：微扰稳定性}

设水银沿管向右移动很小距离 \(x>0\)，此时前端进入水平段但尚未到达
\(B\)。封闭气体在温度 \(T_0\) 下等温膨胀，其压强为

\[
p_g(x)=p_i\frac{q_0}{q_0+x}
=86\frac{90}{90+x}. \tag{4}
\]

前端在水平段内保持与 \(A\) 等高，后端向上移动 \(x\)，两端的竖直高度差由
\(10\,\mathrm{cm}\) 变为 \(10-x\)。若它在新位置静止，需要的气体压强为

\[
p_{\rm eq}(x)=p_0+(10-x)
=86-x, \tag{5}
\]

其中压强和长度均用相应的厘米制数值表示。

在 \(x=0\) 附近展开实际气压：

\[
\begin{aligned}
p_g(x)
&=86\left(1+\frac{x}{90}\right)^{-1}\\
&\approx86\left(1-\frac{x}{90}\right)\\
&=86-\frac{86}{90}x. 
\end{aligned}
\tag{6}
\]

两者之差为

\[
\begin{aligned}
p_g-p_{\rm eq}
&\approx\left(86-\frac{86}{90}x\right)-(86-x)\\
&=\left(1-\frac{86}{90}\right)x\\
&=\frac{2}{45}x>0. 
\end{aligned}
\tag{7}
\]

正微扰后，实际气压大于新位置所需气压，合力仍指向
\(+x\)，使位移继续增大；向左微扰时整柱退回竖直段，高差仍为 \(10\)
cm，压缩气体会产生恢复力；这是弯角处的单侧失稳，不能对正负位移都沿用式（5）。因此

\[
\boxed{\text{初始平衡不稳定}}. \tag{8}
\]

也可比较斜率：\(p_g'(0)=-86/90\)，而
\(p_{\rm eq}'(0)=-1\)。实际气压曲线下降得较慢，故两曲线交点不稳定。

\needspace{3\baselineskip}
\textbf{（三）第（2）问：触发快速流动的温度}

缓慢升温时，在理想的完全无阻力模型中，初态本来就处于不稳定平衡。只要温度比
\(T_0\) 高任意小量，封闭气体在几何位置尚未来得及变化时就有

\[
p_g=p_i\frac{T}{T_0}>p_i, \tag{9}
\]

水银受到向右的净推动，并因不稳定性而加速离开。因此临界值为

\[
\boxed{T_1=T_0=300\,\mathrm K}, \tag{10}
\]

更严格地说，开始运动需要
\(T_1\to300\,\mathrm K^+\)。真实装置若存在静摩擦或黏滞阈值，则需略高温度；题设明确忽略这些阻力，不能人为加入额外温差。

\needspace{3\baselineskip}
\textbf{（四）第（3）问：后端到达 \(A\) 时的速度}

在触发极限 \(T_1\to T_0^+\) 下，可按 \(T_0\)
等温过程计算。水银后端从初始位置到 \(A\)，沿管移动

\[
\Delta s=l_4=10\,\mathrm{cm}. \tag{11}
\]

封闭气柱长度由 \(q_0=90\,\mathrm{cm}\) 增至

\[
q_1=q_0+l_4=100\,\mathrm{cm}. \tag{12}
\]

设管截面积为 \(S\)。封闭气体满足 \(p(q)q=p_iq_0\)，对水银做功

\[
\begin{aligned}
W_g
&=\int_{q_0}^{q_1}p(q)S\,\mathrm dq\\
&=p_iq_0S\int_{q_0}^{q_1}\frac{\mathrm dq}{q}\\
&=p_iq_0S\ln\frac{q_1}{q_0}. 
\end{aligned}
\tag{13}
\]

右端大气对水银做负功

\[
W_0=-p_0S l_4. \tag{14}
\]

还要计算水银质心高度变化。初态水银竖直放置，质心在 \(A\) 下方
\(l_4/2\)。末态水银的前 \(l_2=5\,\mathrm{cm}\) 位于水平段，质心高度为
\(A\) 的高度；后 \(5\,\mathrm{cm}\) 位于右侧竖直段，其质心在 \(A\) 下方
\(2.5\,\mathrm{cm}\)。两半质量相等，末态总质心在 \(A\) 下方

\[
\frac12\times0+\frac12\times\frac{l_4}{4}
=\frac{l_4}{8}. \tag{15}
\]

所以质心上升

\[
\Delta z=\frac{l_4}{2}-\frac{l_4}{8}
=\frac{3l_4}{8}. \tag{16}
\]

水银质量为 \(M_{\rm Hg}=\rho Sl_4\)。由功---能定理

\[
W_g+W_0
=\frac12\rho Sl_4v_1^2
+\rho Sl_4g\frac{3l_4}{8}. \tag{17}
\]

为避免压强单位混乱，把 \(p_i,p_0\)
都写成水银压强头：\(p_i=\rho g(86\,\mathrm{cm})\)，\(p_0=\rho g(76\,\mathrm{cm})\)。式（17）除以
\(\rho Sl_4\)：

\[
\frac{v_1^2}{2g}
=\frac{86\times90}{10}\ln\frac{100}{90}\,\mathrm{cm}
-76\,\mathrm{cm}-\frac{3\times10}{8}\,\mathrm{cm}. \tag{18}
\]

右端压强头数值为

\[
774\ln\frac{10}{9}-76-3.75
=1.799039\,\mathrm{cm}. \tag{19}
\]

取 \(g=9.8\,\mathrm{m/s^2}\)：

\[
v_1=\sqrt{2g\times0.01799039\,\mathrm m}
=0.593811\,\mathrm{m/s}. \tag{20}
\]

故

\[
\boxed{v_1\approx0.594\,\mathrm{m/s}}. \tag{21}
\]

若竞赛中取 \(g=10\,\mathrm{m/s^2}\)，得到约
\(0.600\,\mathrm{m/s}\)，差异只来自 \(g\) 的取值。

\needspace{3\baselineskip}
\textbf{（五）第（4）问：\(B\) 点瞬时压强}

此时气柱长为 \(100\,\mathrm{cm}\)，温度仍为 \(T_0\)，所以水银后端 \(A\)
处的压强为

\[
p_A=p_i\frac{90}{100}
=86\times0.9=77.4\,\mathrm{cmHg}. \tag{22}
\]

水银的一半长度 \(l_2=5\,\mathrm{cm}\) 位于水平段 \(AB\)，另一半位于
\(B\) 以下的竖直段。设沿水银运动方向的加速度大小为 \(a\)。

对水平半柱列牛顿第二定律：

\[
(p_A-p_B)S=\rho Sl_2a. \tag{23}
\]

对竖直向下运动的半柱，以向下为正。上端压强力 \(p_BS\) 向下，下端大气压力
\(p_0S\) 向上，重力 \(\rho Sl_2g\) 向下：

\[
(p_B-p_0)S+\rho Sl_2g
=\rho Sl_2a. \tag{24}
\]

式（23）、式（24）右端相同，故

\[
p_A-p_B=p_B-p_0+\rho gl_2. \tag{25}
\]

解得

\[
p_B=\frac{p_A+p_0-\rho gl_2}{2}. \tag{26}
\]

用厘米水银柱表示：

\[
\boxed{p_B=\frac{77.4+76-5}{2}=74.2\,\mathrm{cmHg}}. \tag{27}
\]

\(p_B\)
可以小于大气压，因为竖直半柱的重力也参与提供向下加速度；这不与右端开口处压强为大气压矛盾。

\needspace{3\baselineskip}
\textbf{（六）第（5）问：能否到达管口}

玻璃管从封闭端到开口的总长度为

\[
L=l_1+l_2+l_3=155\,\mathrm{cm}. \tag{28}
\]

若水银柱前端恰到达开口而尚未流出，封闭气体占据除水银外的全部长度：

\[
q_f=L-l_4=145\,\mathrm{cm}. \tag{29}
\]

从初态到该假想构形，气体做功为

\[
W_g=p_iq_0S\ln\frac{145}{90}. \tag{30}
\]

水银沿管前进距离 \(q_f-q_0=55\,\mathrm{cm}\)，大气做功

\[
W_0=-p_0S(55\,\mathrm{cm}). \tag{31}
\]

初态水银质心在 \(A\) 下方
\(5\,\mathrm{cm}\)；前端到开口时，整柱位于右侧竖直段的最下方，质心在
\(B\) 下方 \(45\,\mathrm{cm}\)。由于 \(A,B\) 等高，质心总共下降
\(40\,\mathrm{cm}\)，重力势能减少，对动能贡献

\[
W_G=\rho Sl_4g(40\,\mathrm{cm}). \tag{32}
\]

把总可用功除以 \(\rho gS\)，其数值为

\[
\begin{aligned}
\frac{W_g+W_0+W_G}{\rho gS}
&=86\times90\ln\frac{145}{90}
-76\times55+10\times40\\
&=-88.6077\,\mathrm{cm^2}<0. 
\end{aligned}
\tag{33}
\]

假想末态若可达，左端应等于非负的动能除以
\(\rho gS\)；现在却为负，说明水银在到达管口前已无足够机械能，必先停下。因此

\[
\boxed{\text{水银不会在这次快速运动中溢出管口}}. \tag{34}
\]

\needspace{3\baselineskip}
\textbf{（七）第（6）问：全部流出的最低温度}

严格无阻力模型在转向点之后会回流，不能永久停在一个新的平衡点。第（6）问另外假定经过很弱的耗散后再次静止；这一耗散只用于建立新的准静态初态，不加入前面快速阶段的能量账本。此后，封闭气体的物质的量没有改变。无论中间平衡位置在哪里，状态常量都可由最初状态写成

\[
\frac{pV}{T}
=\frac{p_i(q_0S)}{T_0}. \tag{35}
\]

继续缓慢升温，要求水银全部流出。临界时考虑``最后一滴水银刚离开管口''的极限：管内剩余水银长度趋于零，静水压头也趋于零；封闭气体充满全管，体积为

\[
V_{\rm last}=LS, \tag{36}
\]

且管口处与大气相通，气体压强趋于

\[
p_{\rm last}=p_0. \tag{37}
\]

代入式（35）：

\[
\frac{p_0LS}{T_{\min}}
=\frac{p_iq_0S}{T_0}. \tag{38}
\]

截面积约去：

\[
T_{\min}=T_0\frac{p_0L}{p_iq_0}. \tag{39}
\]

代入数值：

\[
\begin{aligned}
T_{\min}
&=300\,\mathrm K\times
\frac{76\times155}{86\times90}\\
&=456.589\,\mathrm K. 
\end{aligned}
\tag{40}
\]

故至少需要

\[
\boxed{T_{\min}\approx4.57\times10^2\,\mathrm K}. \tag{41}
\]

临界对象是``最后一滴''，而不是完整 \(10\,\mathrm{cm}\)
水银柱刚到达管口。若仍在末态加入 \(10\,\mathrm{cmHg}\)
的静水压差，就会高估最低温度。

\needspace{3\baselineskip}
\textbf{变式解答：水银全部进入右侧竖管后的平衡}

气柱长度 \(q\) 位于 \(105\le q\le145\) cm
时，水银两端都在右侧竖管，整柱长仍为 10 cm，所需气压恒为

\[
p_{\rm eq}=76-10=66\ {
m cmHg}.
\]

\(T=300\) K 时 \(p_gq=86\times90=7740\)，所以弱阻尼后的稳定平衡为

\[
q_*={7740\over66}\simeq117.2727\ {
m cm}.
\]

此段
\(p_g'(q)<0\)、\(p_{\rm eq}'(q)=0\)：右移扰动使气压低于平衡值，水银被推回；左移扰动则反之。若完全无阻尼，只能讨论围绕该点的振动，不能声称它最后自动静止。

当 \(q\ge145\) cm，水银开始离开出口，剩余柱长为 \(155-q\)，故

\[
p_{\rm eq}=76-(155-q)=q-79.
\]

加热的准静态路径由 \(86\times90(T/300)/q=q-79\) 给出。最后一滴离开时
\(q=155\)、\(p_g=76\)，因此

\[
T_{\rm last}=300{76\times155\over86\times90}\simeq456.589\ {
m K}.
\]

第一滴和最后一滴的条件不同：前者有 10 cm
水银重力压头，后者压头趋零。最后温度不能继续沿用整柱静止的 66 cmHg。

\subsection{六、关键洞察}\label{ux516dux5173ux952eux6d1eux5bdf-6}

\begin{enumerate}
\def\labelenumi{\arabic{enumi}.}
\tightlist
\item
  稳定性由实际气压曲线与平衡所需压强曲线的相对斜率决定，而不只由初态合力为零决定。
\item
  水银跨过弯角时，质心高度必须按各段质量加权；不能只跟踪某一端点。
\item
  \(B\) 点压强是动态压强，必须对两段液柱列动力学方程，不能套静水压公式。
\item
  ``到达管口''和``全部流出''对应不同临界构形；后者由最后一滴极限决定。
\item
  用能量计算出假想末态动能为负，是证明该构形不可达的直接方法。
\end{enumerate}

\subsection{七、易错点}\label{ux4e03ux6613ux9519ux70b9-6}

\begin{enumerate}
\def\labelenumi{\arabic{enumi}.}
\tightlist
\item
  初态只检查 \(p_i=p_0+\rho gl_4\) 就认定稳定。
\item
  把无阻力模型中的触发温度写成明显高于 \(300\,\mathrm K\)。
\item
  气体等温功误写成 \(p\Delta V\)，忽略压强随体积变化。
\item
  漏掉大气压做的负功，或把其位移写成气柱最终长度而非增量。
\item
  计算末态质心时把水平段也计入高度变化。
\item
  用 \(p_B=p_A+\rho gh\)，忽略水银正在加速且 \(AB\) 段水平。
\item
  第（6）问仍把完整水银柱当作末态，未取最后一滴极限。
\item
  将 \(\mathrm{cmHg}\) 与帕斯卡直接混入同一式子，造成量纲错误。
\end{enumerate}

\subsection{八、课堂讲解建议}\label{ux516bux8bfeux5802ux8bb2ux89e3ux5efaux8bae-6}

先画 \(p_g(x)\) 与 \(p_{\rm eq}(x)\)
两条曲线，只讲初态附近斜率。学生理解``离开后净力同向''后，再进入后续能量计算，整道题的运动方向会更清楚。

第（3）问建议让学生分别报出气体功、大气功和质心高度变化，三项全部确认后再合并。可用``单位质量、单位重力对应多少厘米压强头''的方式统一单位，避免在厘米水银柱和
SI 制之间反复换算。

第（4）问必须在题图上把水银涂成水平、竖直各
\(5\,\mathrm{cm}\)，并给两个半柱标同一个沿管加速度
\(a\)。第（5）、（6）问则各画一个末态图，对比``前端刚到口''和``最后一滴刚离口''。

备课时可把第（3）问的能量账本按``施功者---受功者---正负号---几何位移''逐项核对。封闭气体对水银做正功，右端大气做负功，质心上升使一部分净功转化为重力势能，剩余部分才是动能。若学生算出的
\(v_1\) 明显大于
\(1\,\mathrm{m/s}\)，通常不是开方问题，而是漏除了水银质量中的长度
\(l_4\)，或漏记大气功。

还可让学生对式（18）做一次独立量纲审查：等号右侧三项都必须是长度，左侧
\(v^2/(2g)\) 也是长度。只有在这一步确认后，才能把
\(1.799039\,\mathrm{cm}\) 换成
\(0.01799039\,\mathrm m\)。这种``先保留压强头、最后一次换 SI
单位''的策略适用于几乎所有细管液柱题。

第（5）问的负值不是``负动能''，而是反证假想末态不可达。讲解时应明确：能量守恒式只对真实可达状态成立；把管口构形当作试探终点后，计算得到需要额外输入
\(88.6077\,\rho gS\,\mathrm{cm^2}\)
的能量，所以实际轨迹必在此前出现速度为零的转折点。

\subsection{九、变式训练}\label{ux4e5dux53d8ux5f0fux8badux7ec3-6}

\begin{enumerate}
\def\labelenumi{\arabic{enumi}.}
\tightlist
\item
  若初始气柱长度为一般值 \(q_0\)，推导平衡稳定的条件。
\item
  若管壁绝热，快速阶段改用
  \(pV^\gamma=\text{常量}\)，重写第（3）问的气体功。
\item
  改变水平段长度 \(l_2\)，求水银后端到达 \(A\) 时的质心高度和 \(B\)
  点压强。
\item
  求水银第一次停下时的气柱长度，并讨论该新平衡的稳定性。
\item
  若存在恒定阻力 \(F_f\)，写出第（5）问的可达判据。
\end{enumerate}

\needspace{3\baselineskip}
\textbf{用同一个气柱长度写出五段平衡压力}

取封闭端至水银后端的气柱长度为 \(q\)。原初值是 \(q=90\) cm，后端到弯角
\(A\) 还差 10 cm。由于横管长只有5 cm，先后出现五个几何分段：

\[
p_{\rm eq}(q)=\begin{cases}
176-q,&90\le q\le95,\\
271-2q,&95\le q\le100,\\
171-q,&100\le q\le105,\\
66,&105\le q\le145,\\
q-79,&145\le q\le155.
\end{cases}
\]

各段接点的压力分别为 \(86,81,71,66,66,76\) cmHg，必须连续。第二段斜率为
\(-2\)，因为一端在左竖管、另一端已进右竖管，两端高度变化共同改变压头；其他进入段只有一个端点贡献变化。

水银未流出时，气体做功减去平衡压力功就是动能增加：

\[
\Delta K=\rho gS\int_{q_0}^{q}[p_g(\tilde q)-p_{\rm eq}(\tilde q)]\,\mathrm d\tilde q.
\]

压力和长度均以cm为单位计算时，积分单位为cm²，最后换成SI。\(q=100\)
时把前两段分别积分，得到 \(\Delta K/(\rho gS)=17.99039\)
cm²；再除原水银柱长10 cm，可得 \(v^2/(2g)=1.799039\) cm。

初态右移的恢复性由两曲线斜率比较，不能由某个终态速度是否实数倒推。图中的橙线必须经过全部接点，尤其不能省略长度只有5
cm的第二段。

\newpage

\section{第 8 题}\label{ux7b2c-8-ux9898}

在实验室参考系中，一个正电子静止，静质量为
\(m\)；一个电子沿直线飞来，其总能量为 \(\gamma mc^2\)，其中
\(\gamma\geq1\)。二者碰撞后湮灭，产生频率分别为 \(\nu_1,\nu_2\)
的两个光子。两光子在实验室系中的夹角为 \(\theta\)，普朗克常量为 \(h\)。

（1）当 \(\nu_1,\nu_2\) 满足什么关系时，夹角 \(\theta\) 取得最小值？

（2）求最小夹角 \(\theta_{\min}(\gamma)\)。

\begin{figure}[H]
\centering
\includegraphics[width=0.82\linewidth,height=\textheight,keepaspectratio,alt={图8-1　实验室系中电子---正电子湮灭为两个光子的动量几何}]{../figures/question-08.pdf}
\caption*{图8-1　实验室系中电子---正电子湮灭为两个光子的动量几何}
\end{figure}

\subsection{一、题型定位}\label{ux4e00ux9898ux578bux5b9aux4f4d-7}

两光子必须同时满足能量和动量守恒。利用总四动量不变量可把夹角写成两个光子能量的函数；在总能量固定时，等能分配给出最小夹角。

\subsection{二、知识点}\label{ux4e8cux77e5ux8bc6ux70b9-7}

\begin{enumerate}
\def\labelenumi{\arabic{enumi}.}
\tightlist
\item
  相对论能量---动量关系 \(E^2=p^2c^2+m^2c^4\)。
\item
  系统总四动量平方的洛伦兹不变性。
\item
  光子能量与动量 \(E_\gamma=pc=h\nu\)。
\item
  两向量和的余弦定理。
\item
  定和条件下乘积的最大值。
\item
  质心速度与洛伦兹变换的几何解释。
\item
  非相对论极限与超相对论极限检查。
\end{enumerate}

\subsection{三、物理图像}\label{ux4e09ux7269ux7406ux56feux50cf-7}

初态总动量沿入射电子方向。末态两光子的动量矢量首尾相加必须等于该总动量，能量之和则等于
\((\gamma+1)mc^2\)。当一个光子分走过多能量时，两动量矢量长度相差很大，为得到固定的合动量，它们之间需要更大的张角；当两者等能时，动量三角形关于总动量方向对称，夹角最小。

在质心系中，总动量为零，所以两个光子的动量大小相等、方向相反。把这一构形沿初态总动量方向推进到实验室系，两个光子的纵向动量同时增加。若它们在质心系中垂直于推进方向发射，变换后在实验室系中仍等能，并形成最小夹角。

\subsection{四、解题入口}\label{ux56dbux89e3ux9898ux5165ux53e3-7}

避免分别列 \(x,y\) 两方向的多个三角方程。先求初态系统的

\[
\mathcal S=E_{\rm tot}^2-p_{\rm tot}^2c^2,
\]

再在末态把同一不变量写成

\[
2E_1E_2(1-\cos\theta).
\]

能量守恒固定 \(E_1+E_2\)。由于 \(0\leq\theta\leq\pi\)
上余弦单调递减，最小化 \(\theta\) 等价于最大化
\(\cos\theta\)，继而等价于最大化 \(E_1E_2\)。

\subsection{五、解析}\label{ux4e94ux89e3ux6790-7}

\begin{figure}[H]
\centering
\includegraphics[width=0.855\linewidth,height=\textheight,keepaspectratio,alt={图8-2　等能极值构形与质心系反向双光子图像}]{../figures/solution-08.pdf}
\caption*{图8-2　等能极值构形与质心系反向双光子图像}
\end{figure}

图8-2中，左图给出一个实际满足动量合成的等能例子，右图给出对应的质心系反向光子。由不变量固定
\(E_1E_2(1-\cos\theta)\)，即可证明等能时夹角最小。

\needspace{3\baselineskip}
\textbf{（一）初态总能量和总动量}

静止正电子的能量为

\[
E_+=mc^2, \tag{1}
\]

动量为零。入射电子的总能量为

\[
E_-=\gamma mc^2. \tag{2}
\]

由单粒子能量---动量关系

\[
E_-^2=p^2c^2+m^2c^4, \tag{3}
\]

得到电子动量大小

\[
\begin{aligned}
p^2c^2
&=\gamma^2m^2c^4-m^2c^4\\
&=(\gamma^2-1)m^2c^4, 
\end{aligned}
\tag{4}
\]

即

\[
\boxed{pc=mc^2\sqrt{\gamma^2-1}}. \tag{5}
\]

因此初态系统总能量和总动量大小分别为

\[
\boxed{E=(\gamma+1)mc^2},\qquad
\boxed{p_{\rm tot}=p}. \tag{6}
\]

\needspace{3\baselineskip}
\textbf{（二）计算系统四动量不变量}

初态系统总四动量平方为

\[
\begin{aligned}
\mathcal S
&=E^2-p_{\rm tot}^2c^2\\
&=(\gamma+1)^2m^2c^4-(\gamma^2-1)m^2c^4\\
&=[\gamma^2+2\gamma+1-\gamma^2+1]m^2c^4\\
&=2(\gamma+1)m^2c^4. 
\end{aligned}
\tag{7}
\]

所以

\[
\boxed{\mathcal S=2(\gamma+1)m^2c^4}. \tag{8}
\]

这个量在所有惯性系中相同，也是末态双光子系统必须保持的约束。

\needspace{3\baselineskip}
\textbf{（三）末态双光子的能量与夹角关系}

记

\[
E_1=h\nu_1,\qquad E_2=h\nu_2. \tag{9}
\]

光子动量大小为

\[
p_1=\frac{E_1}{c},\qquad p_2=\frac{E_2}{c}. \tag{10}
\]

能量守恒给出

\[
E_1+E_2=E=(\gamma+1)mc^2. \tag{11}
\]

动量守恒要求 \(\bm p_1+\bm p_2=\bm p_{\rm tot}\)。对两边平方：

\[
p_{\rm tot}^2=p_1^2+p_2^2+2p_1p_2\cos\theta. \tag{12}
\]

乘以 \(c^2\)：

\[
p_{\rm tot}^2c^2
=E_1^2+E_2^2+2E_1E_2\cos\theta. \tag{13}
\]

另一方面

\[
E^2=(E_1+E_2)^2
=E_1^2+E_2^2+2E_1E_2. \tag{14}
\]

式（14）减式（13）：

\[
\begin{aligned}
E^2-p_{\rm tot}^2c^2
&=2E_1E_2(1-\cos\theta). 
\end{aligned}
\tag{15}
\]

代入不变量式（8）：

\[
1-\cos\theta
=\frac{(\gamma+1)m^2c^4}{E_1E_2}. \tag{16}
\]

或写为

\[
\boxed{
\cos\theta
=1-\frac{\mathcal S}{2E_1E_2}}. \tag{17}
\]

\needspace{3\baselineskip}
\textbf{（四）第（1）问：频率关系的极值判定}

式（11）固定了 \(E_1+E_2\)，而 \(\mathcal S\) 也是固定量。在
\(0\leq\theta\leq\pi\) 内，要使 \(\theta\) 最小，应使 \(\cos\theta\)
最大。由式（17），这要求 \(E_1E_2\) 最大。

利用

\[
(E_1-E_2)^2\geq0, \tag{18}
\]

展开得

\[
(E_1+E_2)^2\geq4E_1E_2. \tag{19}
\]

所以在定和条件下

\[
E_1E_2\leq\frac{(E_1+E_2)^2}{4}, \tag{20}
\]

等号当且仅当

\[
E_1=E_2. \tag{21}
\]

由于 \(E_i=h\nu_i\)，最小夹角条件为

\[
\boxed{\nu_1=\nu_2}. \tag{22}
\]

此时每个光子的能量和频率为

\[
E_1=E_2=\frac{\gamma+1}{2}mc^2,\qquad
\nu_1=\nu_2=\frac{(\gamma+1)mc^2}{2h}. \tag{23}
\]

\needspace{3\baselineskip}
\textbf{（五）第（2）问：最小夹角}

把 \(E_1E_2=E^2/4\) 代入式（17）：

\[
\begin{aligned}
\cos\theta_{\min}
&=1-\frac{\mathcal S}{2(E^2/4)}\\
&=1-\frac{2\mathcal S}{E^2}\\
&=1-\frac{2\cdot2(\gamma+1)m^2c^4}
{(\gamma+1)^2m^2c^4}\\
&=1-\frac4{\gamma+1}\\
&=\frac{\gamma-3}{\gamma+1}. 
\end{aligned}
\tag{24}
\]

因此

\[
\boxed{
\theta_{\min}(\gamma)
=\arccos\frac{\gamma-3}{\gamma+1}}. \tag{25}
\]

用半角公式

\[
\cos^2\frac\theta2=\frac{1+\cos\theta}{2}, \tag{26}
\]

可得更有物理意义的形式：

\[
\begin{aligned}
\cos^2\frac{\theta_{\min}}2
&=\frac12\left(1+\frac{\gamma-3}{\gamma+1}\right)\\
&=\frac{\gamma-1}{\gamma+1}. 
\end{aligned}
\tag{27}
\]

因为 \(0\leq\theta_{\min}/2\leq\pi/2\)，取非负平方根：

\[
\boxed{
\theta_{\min}(\gamma)
=2\arccos\sqrt{\frac{\gamma-1}{\gamma+1}}}. \tag{28}
\]

式（25）与式（28）完全等价。

\needspace{3\baselineskip}
\textbf{（六）质心系解释}

初态系统质心相对实验室系的速度为

\[
\beta_c\equiv\frac{V_c}{c}
=\frac{p_{\rm tot}c}{E}. \tag{29}
\]

代入式（5）、式（6）：

\[
\begin{aligned}
\beta_c
&=\frac{\sqrt{\gamma^2-1}}{\gamma+1}\\
&=\sqrt{\frac{(\gamma-1)(\gamma+1)}{(\gamma+1)^2}}\\
&=\boxed{\sqrt{\frac{\gamma-1}{\gamma+1}}}. 
\end{aligned}
\tag{30}
\]

比较式（27）和式（30）：

\[
\boxed{\cos\frac{\theta_{\min}}2=\beta_c}. \tag{31}
\]

在质心系中两光子等能、反向。若它们相对于质心系推进方向横向发射，洛伦兹变换到实验室系后，两光子纵向动量相等，横向动量相反，因此实验室系中仍然等能并关于总动量方向对称。每条光子与总动量方向夹角为
\(\theta_{\min}/2\)，其方向余弦正是
\(\beta_c\)。这给出了式（31）的几何解释。

\needspace{3\baselineskip}
\textbf{（七）极限与反代检查}

当 \(\gamma=1\) 时，电子也静止，初态总动量为零。式（24）给出

\[
\cos\theta_{\min}=-1,\qquad
\theta_{\min}=\pi, \tag{32}
\]

即两光子反向，符合动量守恒。

当 \(\gamma=3\) 时，

\[
\cos\theta_{\min}=0,\qquad
\theta_{\min}=\frac\pi2. \tag{33}
\]

当 \(\gamma\to\infty\) 时，\(\beta_c\to1\)，故
\(\theta_{\min}\to0\)。由式（28）还可得到渐近形式

\[
\theta_{\min}\sim2\sqrt{\frac{2}{\gamma+1}}, \tag{34}
\]

说明高能极限下两个光子在实验室系中被束射到入射方向附近。

\subsection{六、关键洞察}\label{ux516dux5173ux952eux6d1eux5bdf-7}

\begin{enumerate}
\def\labelenumi{\arabic{enumi}.}
\tightlist
\item
  双光子末态的夹角受系统四动量不变量控制，而不是只由普通能量守恒决定。
\item
  总能量固定时，等能分配使 \(E_1E_2\) 最大，从而使动量三角形的张角最小。
\item
  实验室系的等能不意味着任意质心系发射方向；极值构形对应质心系横向发射。
\item
  半角余弦等于质心速度的无量纲值，使代数答案具有清晰的相对论几何意义。
\end{enumerate}

\subsection{七、易错点}\label{ux4e03ux6613ux9519ux70b9-7}

\begin{enumerate}
\def\labelenumi{\arabic{enumi}.}
\tightlist
\item
  把入射电子动量误写成 \(\gamma mvc\) 或直接写成
  \(\gamma mc\)，没有使用能量---动量关系。
\item
  认为湮灭产生的两个光子必然反向；这只在总动量为零的参考系成立。
\item
  从式（17）判断极值时方向颠倒，错误地最小化 \(E_1E_2\)。
\item
  漏掉静止正电子的静能 \(mc^2\)。
\item
  把 \(E_1+E_2\) 固定误解为 \(\nu_1,\nu_2\)
  可取零；动量约束还要求式（17）的余弦在 \([-1,1]\) 内。
\item
  半角开平方时随意取负号，忽略夹角范围。
\end{enumerate}

\subsection{八、课堂讲解建议}\label{ux516bux8bfeux5802ux8bb2ux89e3ux5efaux8bae-7}

先画末态动量三角形，让学生说出三条边分别是
\(E_1/c,E_2/c,p_{\rm tot}\)。随后用余弦定理与能量定和直接观察：当两条光子动量边相等时图形最对称，张角最小。

四动量方法讲完后，再用质心系做第二次解释。两种方法承担不同功能：不变量法负责严密计算，质心系负责建立直觉。最后用
\(\gamma=1,3,\infty\) 三个点检查公式，帮助学生把表达式与图形联系起来。

板书时建议把``初态不变量''\,``末态不变量''\,``定和极值''分成连续三栏。第一栏只出现
\(\gamma,m,c\)；第二栏只出现 \(E_1,E_2,\theta\)；第三栏再用
\(E_1+E_2=E\)
连接二者。这样学生能看清不变量承担的是跨越反应前后的桥梁，而不是又一条独立的能量守恒式。

频率相等的结论应在求角度之前得到，并用动量图作方向检查。等长的两光子动量矢量关于合动量方向对称，夹角的一半就是每个光子偏离总动量方向的角度。若某份计算在
\(\gamma=1\) 时仍给出小于 \(\pi\)
的夹角，说明漏掉了静止正电子的静能或不变量中少了交叉项。

对基础较弱的学生，可先不用四矢量术语，只把 \(E^2-p^2c^2\)
作为由初态直接算出、在反应前后保持不变的组合；等学生完成式（15）后，再指出这正是总四动量的平方。这样既不牺牲严密性，也能避免概念名称先于计算对象出现。

\subsection{九、变式训练}\label{ux4e5dux53d8ux5f0fux8badux7ec3-7}

\begin{enumerate}
\def\labelenumi{\arabic{enumi}.}
\tightlist
\item
  已知两光子能量比 \(E_1/E_2=\eta\)，求实验室系夹角
  \(\theta(\eta,\gamma)\)。
\item
  若入射粒子与静止靶粒子的静质量不同，建立对应的不变量和最小夹角条件。
\item
  在质心系中设光子与推进方向夹角为
  \(\chi\)，经洛伦兹变换求实验室系两光子能量比。
\item
  证明在固定初态四动量下，等能双光子构形确实给出全局最小而非局部极值。
\item
  令两光子频率存在小偏差 \(\nu_{1,2}=\bar\nu(1\pm\varepsilon)\)，展开到
  \(\varepsilon^2\)，求夹角相对最小值的增量。这一变式用于说明等能点的一阶变化为零、极小值由二阶项控制。
\item
  给定实验室系测得的 \(\theta\) 与总光子能量 \(E_1+E_2\)，反求入射电子的
  \(\gamma\)，并讨论只有角度而没有能量标定时能否唯一确定入射能量。
\item
  若只探测到一个光子的频率与方向，利用能量---动量守恒建立另一个光子的能量和方向方程，说明哪些测量量足以重建初态。
\item
  从质心系出发写出光子能量的多普勒变换 \[
  E_{1,2}=\Gamma_cE'\bigl(1\pm\beta_c\cos\chi\bigr),
  \] 验证 \(E_1=E_2\) 要求 \(\cos\chi=0\)，并由动量分量再次推出
  \(\cos(\theta_{\min}/2)=\beta_c\)。
\end{enumerate}

\needspace{3\baselineskip}
\textbf{变式解答：固定能量不等分时的角度}

令两个光子能量分别为 \(fE\)、\((1-f)E\)，\(0<f<1\)。原总能量
\(E=(\gamma+1)mc^2\)，不变量给出

\[
1-\cos\theta={1\over(\gamma+1)f(1-f)}.
\]

由于 \(f(1-f)\le1/4\)，右端在 \(f=1/2\) 时最小，也即夹角最小。并非所有
\(f\) 都允许：还须满足 \(1-\cos\theta\le2\)，所以

\[
f(1-f)\ge{1\over2(\gamma+1)}.
\]

例如 \(\gamma=3\)，等能时 \(\theta=90^\circ\)；取 \(f=1/3\) 则

\[
\cos\theta=1-{1\over4(1/3)(2/3)}=-{1\over8},
\]

夹角约 \(97.18^\circ\)，确实增大。可行能量份额为
\(\tfrac12(1\pm1/\sqrt2)\)
之间。端点对应两光子反向，其能量差承担非零总动量。

等能最小角还可从图中对称动量三角形直接核对：每个光子沿入射方向的投影合计
\(E\cos(\theta/2)/c\)，须等于初电子动量
\(mc\sqrt{\gamma^2-1}\)。这是独立于不变量展开的第二条检查。
\end{document}
